% !TeX root = ../../Book2.tex
% 纲要标题：## 第1章　线性代数的结构观点
% 纲要位置：计划纲要.md，第 829 行。

\chapter{线性代数的结构观点}
\label{b2:ch:01}
\BookChapterNumber{b2:ch:01}

\begin{chapterabstract}
本章重看向量空间中不依赖坐标的结构：核、像、商空间、对偶映射与双对偶求值映射，并分清有限维和任意维结论的适用范围。线性算子还被解释为多项式环的作用，由此提出循环子空间和标准形的分类问题；有理标准形与 Jordan 标准形的证明见第四章。
\end{chapterabstract}

\begin{learninggoals}
\begin{itemize}
\item 用基扩充、商空间及核与像构造线性映射，说明补空间的选择与商空间本身的区别。
\item 写出对偶映射和双对偶求值映射，判断哪些结论需要有限维，并用无限坐标空间给出反例。
\item 计算循环子空间和最小多项式，理解不变子空间与多项式作用之间的关系。
\item 区分矩阵等价与相似，检验常见不变量，并辨认标准形中尚需证明的存在性和唯一性。
\end{itemize}
\end{learninggoals}

线性方程组可以用矩阵求解，但一张矩阵并不能说明解答中哪些选择是必要的。例如投影 \(T:k^2\to k^2\)，\(T(x,y)=(x,0)\)，其核是第二坐标轴。由 \(T\) 直接得到商映射及商到像的映射
\[
q(x,y)=(x,y)+\ker T,\qquad
\bar T\bigl((x,y)+\ker T\bigr)=(x,0).
\]
同一商类中的向量具有相同的第一坐标，所以 \(\bar T\) 与代表元无关；它把商空间同构到第一坐标轴，无须选择补空间。若进一步要在 \(k^2\) 中为每个商类选一个代表，则可以选 \((x,0)\)，也可以选 \((x,x)\)，分别落在第一坐标轴和直线 \(y=x\) 上。这两种线性选取都满足 \(q(r(q(x,y)))=q(x,y)\)，但对应不同的映射 \(r:k^2/\ker T\to k^2\)。依赖补空间的是这一步挑代表元，而不是商到像的同构。

\ProseParagraphBreak

这个区别不只发生在商空间。选基把向量变成坐标列，因而便于计算；对偶映射、求值映射和线性算子的复合却应当在换基之后仍表达同一件事。我们将从熟悉的消元和维数公式出发，逐步辨认哪些同构由定义直接给出，哪些需要任意选基或补空间。到讨论单个算子时，还要把它对向量的反复作用合在一起，才能提出不依赖矩阵写法的分类问题。

\ProseParagraphBreak

本章以线性方程组、基本矩阵运算、行列式，以及域与一元多项式的带余除法为基础。这里会重新写出向量空间及其映射的定义，补齐后来模论实际使用的构造。任意维基的存在沿用第一卷附录 A 的选择公理；有限维的基扩充则只需有限步骤。无限维对偶的例子用于说明有限维结论的局限，不涉及拓扑向量空间。

\ProseParagraphBreak
向量空间与线性映射可参阅 Roman 的《Advanced Linear Algebra》\cite[Chapters 1--2]{Roman2008AdvancedLinearAlgebra}；商空间与对偶理论见同书\cite[Chapter 3, pp. 87--105]{Roman2008AdvancedLinearAlgebra}。Shoup 的《A Computational Introduction to Number Theory and Algebra》\cite[\S\S 13.4--13.5, 14.2]{Shoup2008NumberTheoryAlgebra}适合结合有限域上的计算复习基、维数及矩阵表示。Shoup 在第 14.2 节使用行向量右乘矩阵的约定；与本章的列向量约定对照时，须同时转换矩阵表示与复合公式。算子的模论解释还可对照 Roman 同书\cite[Chapter 7, The Module Associated with a Linear Operator, pp. 164--167]{Roman2008AdvancedLinearAlgebra}，其中所用的一般模概念将在下一章定义。

% !TeX root = ../../Book2.tex
% 纲要标题：### 1.1 向量空间与映射
% 纲要位置：计划纲要.md，第 831 行。

\section{向量空间与映射}
\label{b2:sec:0101}

给定线性映射 \(f:V\to W\)，两个向量何时被它送到同一点？把 \(\ker f\) 中的向量视为零后，为什么所得商空间 \(V/\ker f\) 恰好同构于 \(\operatorname{im}f\)？本节从向量空间的基和商空间出发回答这些问题，并说明补空间虽可用来选取商类的代表，却依赖选择。有限维部分的证明只涉及有限线性组合；任意维部分则明确说明基的存在所需的选择原理。

% 纲要内容（仅作写作计划，尚非教材正文）：
%
% * 有限维和任意维向量空间的约定
% * 基、维数与商空间；任意维数的基数加法公式与补空间到像的同构
% * 线性映射的核、像与秩
%

\subsection{有限维与任意维向量空间}
\label{b2:subsec:0101:01}

本章的 \(k\) 总表示交换域，\(0\ne1\)；向量空间中的标量默认写在左侧。

\begin{definition}\label{b2:def:vector-space}
设 \(k\) 是域，\((V,+)\) 是交换群，并给定标量乘法 \(k\times V\to V\)。若对任意 \(a,b\in k\) 及 \(u,v\in V\)，下列等式均成立：
\[
a(u+v)=au+av,\qquad (a+b)v=av+bv,\qquad
(ab)v=a(bv),\qquad 1v=v.
\]
满足这些条件的 \(V\) 称为 \(k\) 上的\term[xiangliang-kongjian]{向量空间}[vector space]。包含 \(0\) 且对加法与标量乘法封闭的子集 \(U\subseteq V\)，称为 \(V\) 的\term[zi-kongjian]{子空间}[subspace]。
\end{definition}

标量的零和向量的零均写作 \(0\)，由所在等式区分。公理给出 \(0v=0\)、\(a0=0\) 及 \((-1)v=-v\)；例如 \(0v=(0+0)v=0v+0v\)，在加法群中消去 \(0v\) 即得第一式。子空间因而也对取负封闭。零空间 \(\{0\}\) 是向量空间，其基将在下面约定为空集。

\ProseParagraphBreak
坐标空间 \(k^n\) 和多项式空间 \(k[t]\) 都按分量或系数运算。还要区分两种无限坐标空间：给定集合 \(I\)，记
\[
k^I=\{(a_i)_{i\in I}\mid a_i\in k\},\qquad
k^{(I)}=\{(a_i)_{i\in I}\in k^I\mid a_i\ne0\;\text{的指标只有有限个}\}.
\]
二者都是向量空间，后者是前者的子空间。\term[zhiji]{支集}[support] \(\operatorname{supp}(a)=\{i\in I\mid a_i\ne0\}\) 记录非零坐标。若 \(I\) 无限，常数族 \((1)_{i\in I}\) 属于 \(k^I\)，却不属于 \(k^{(I)}\)。这里的向量空间只使用有限和；没有给定拓扑时，\(\sum_{i\in I}a_iv_i\) 不能任意解释为无限和。
\symidx{\(k^I,k^{(I)}\)：全部标量族及有限支集标量族}

\begin{definition}\label{b2:def:linear-span-basis}
设 \(V\) 是 \(k\) 向量空间，\(S\subseteq V\)。\(S\) 中向量的全部有限线性组合构成\term[xianxing-shengcheng-kongjian]{线性生成空间}[linear span] \(\operatorname{span}_kS\)；约定空和为 \(0\)。线性生成空间也称线性包络。\footnote{“线性生成空间”见黎景辉等《高等线性代数学》\cite[定义 1.13, p. 4]{LiBaiZhou2014AdvancedLinearAlgebra}；“线性包络”见席南华《基础代数》第二卷\cite[§1.2, p. 5]{Xi2018BasicAlgebraII}。}\termidx[xianxing-baoluo]{线性包络}若 \(\operatorname{span}_kS=V\)，就称 \(S\) 生成 \(V\)。给定指标集 \(I\) 上的向量族 \((v_i)_{i\in I}\)，如果每个有限关系
\[
\sum_{i\in F}a_iv_i=0\quad(F\subseteq I\;\text{有限})
\]
（其中 \(a_i\in k\)）都迫使所有 \(a_i=0\)，就称该向量族\term[xianxing-wuguan]{线性无关}[linearly independent]。如果它还生成 \(V\)，就称它为 \(V\) 的一组\term[ji]{基}[basis]。如果 \(V\) 有有限生成集，就称它为\term[youxian-wei]{有限维}[finite-dimensional]空间；否则称为无限维空间。
\end{definition}
\symidx{\(\operatorname{span}_kS\)：向量集的线性生成空间}

线性生成空间包含 \(S\)，对加法和标量乘法封闭，并包含于每个含 \(S\) 的子空间，故是含 \(S\) 的最小子空间。基的两个条件合起来，恰好说明每个向量都能唯一地写成基向量的有限线性组合。坐标向量 \((e_i)_{i\in I}\) 是 \(k^{(I)}\) 的基；\(1,t,t^2,\ldots\) 是 \(k[t]\) 的基。\(I\) 无限时，同一族 \((e_i)\) 不能生成 \(k^I\)，所以不是后者的基。

\subsection{基、维数与商空间}
\label{b2:subsec:0101:02}

在有限生成的空间中，可以逐个删去多余的生成元，直到得到一组基。比较两组基时，需要下面的交换论证。

\begin{lemma}[有限交换引理]\label{b2:lem:finite-exchange}
设 \(V\) 是 \(k\) 向量空间。若 \(v_1,\ldots,v_n\) 生成 \(V\)，而 \(u_1,\ldots,u_m\in V\) 线性无关，则 \(m\le n\)，且可用这些 \(u_i\) 替换生成组中的 \(m\) 个向量，得到仍生成 \(V\) 的 \(n\) 个向量。
\end{lemma}
\begin{proof}
把 \(u_1\) 写成 \(v_j\) 的线性组合。因为 \(u_1\ne0\)，某个系数非零；解出对应的 \(v_j\)，即可用 \(u_1\) 替换它而不改变生成空间。

假设已用 \(u_1,\ldots,u_r\) 替换了 \(r\) 项。把 \(u_{r+1}\) 写成当前生成组的线性组合；至少一个尚未替换的 \(v_j\) 的系数非零，否则 \(u_{r+1}\) 属于前 \(r\) 个 \(u_i\) 的线性生成空间，违反无关性。解出这个 \(v_j\) 并替换，便完成下一步。若已替换全部 \(n\) 项，就不能再加入一个与前面无关的 \(u_i\)，所以 \(m\le n\)。
\end{proof}

\begin{theorem}[基的扩充]\label{b2:thm:vector-basis-extension}
设 \(V\) 是 \(k\) 向量空间。\(V\) 中的每个线性无关子集都可扩充为一组基。若 \(V\) 由有限集 \(S\) 生成，可以只从 \(S\) 中选取所需的新向量。
\end{theorem}
\begin{proof}
先看有限生成情形。由\cref{b2:lem:finite-exchange}，给定无关集不可能有超过 \(|S|\) 个元素，否则取出 \(|S|+1\) 个就得到矛盾。将 \(S\) 的元素依次加入：已在当前线性生成空间中的跳过，不在其中的加入。后一操作保持无关性，因为一个非平凡新关系若含新向量的非零系数，就会把它写进原来的线性生成空间。有限步后得到含原无关集、且线性生成空间包含 \(S\) 的基。

任意维情形采用第一卷附录A中的基扩充命题。该命题用 Zorn 引理\footnote{佐恩（Max August Zorn，1906-1993），德国裔美国数学家，提出与选择公理等价的极大原则，并研究非结合代数。}考察包含给定无关集的全部无关子集：链的并仍无关，因为每个关系只涉及有限多个向量；极大无关集若不生成全空间，还能加入线性生成空间以外的向量，矛盾。这里沿用第一卷已经采用的选择公理，不把无限扩充理解为一个必会在有限步停止的算法。
\end{proof}

\begin{theorem}[维数不变性]\label{b2:thm:vector-dimension}
设 \(V\) 是 \(k\) 向量空间。\(V\) 的任意两组基等势；其共同基数称为 \(V\) 的\term[weishu]{维数}[dimension]，记作 \(\dim_kV\)。特别地，零空间的维数为 \(\dim_k0=0\)。
\end{theorem}
\symidx{\(\dim_kV\)：向量空间的维数}
\begin{proof}
若一组基 \(B\) 有限，\cref{b2:lem:finite-exchange}说明另一组基 \(C\) 至多有 \(|B|\) 个元素，反过来又有 \(|B|\le|C|\)，故大小相同。

若两组基都无限，把每个 \(c\in C\) 写为 \(B\) 的有限线性组合。这些支集的并覆盖 \(B\)：否则一个从未出现的 \(b\in B\) 就不能由 \(C\) 生成。第一卷附录A中的有限支集计数命题说明，由 \(|C|\) 个有限集合组成的并至多有基数 \(|C|\)，所以 \(|B|\le|C|\)。交换两组基得到反向不等式，再由第一卷附录A中的 Cantor--Bernstein 定理\footnote{康托尔（Georg Cantor，1845-1918），德国数学家，创立集合论并研究无限基数。伯恩施坦（Felix Bernstein，1878-1956），德国数学家，在集合论中证明了这一等势判据。}得到等势。
\end{proof}

以下维数都指基的基数；有限维时它是非负整数，任意维时相加则指不交并所给的基数加法。

\begin{definition}\label{b2:def:quotient-vector-space}
设 \(V\) 是 \(k\) 向量空间，\(U\le V\) 是子空间。对 \(v,w\in V\) 和 \(a\in k\)，在加法商群 \(V/U\) 上定义
\[
(v+U)+(w+U)=(v+w)+U,\qquad a(v+U)=av+U.
\]
所得向量空间称为\term[shang-kongjian]{商空间}[quotient space]；自然映射 \(q:V\to V/U\) 取 \(q(v)=v+U\)。
\end{definition}
\symidx{\(V/U\)：向量空间的商空间}

若更换 \(v,w\) 的代表元，变化量都在 \(U\) 中，相加或乘标量后仍在 \(U\)，所以两种运算良定义。向量空间公理逐项从 \(V\) 下降，例如 \(a(b(v+U))=(ab)v+U\)。两个向量在商中相同，当且仅当它们的差在 \(U\) 中；商空间将 \(U\) 内的差别视为零。

\begin{proposition}\label{b2:prop:quotient-basis}
设 \(V\) 是 \(k\) 向量空间，\(U\le V\)。把 \(U\) 的一组基 \(B_U\) 扩充为 \(V\) 的基 \(B_U\sqcup C\)，则 \((c+U)_{c\in C}\) 是 \(V/U\) 的基，且有基数加法公式
\[
\dim_kV=\dim_kU+\dim_k(V/U).
\]
令 \(W=\operatorname{span}_kC\)，每个 \(v\in V\) 唯一写为 \(u+w\)，其中 \(u\in U,w\in W\)。若一个子空间 \(W\le V\) 满足这样的唯一分解，就记 \(V=U\oplus W\)，并称 \(W\) 为 \(U\) 的一个\term[bu-kongjian]{补空间}[complement]。
\end{proposition}
\begin{proof}
基展开保证 \(v=u+\sum_{c\in F}a_cc\)，所以这些陪集生成商空间。若 \(\sum a_c(c+U)=0\)，则 \(\sum a_cc\in U\)，可用 \(B_U\) 表示。移到等号一侧，再用 \(B_U\sqcup C\) 的无关性，得到全部 \(a_c=0\)。因此 \(V/U\) 的基与 \(C\) 等势，而 \(V\) 的基是不交并 \(B_U\sqcup C\)，给出所述维数公式。基展开也给出 \(V=U+W\)，且刚才的无关性说明 \(U\cap W=0\)，从而 \(u+w\) 的分解唯一。
\end{proof}

补空间依赖选择。例如在 \(k^2\) 中，\(U=ke_1\) 的补空间可以取 \(ke_2\)，也可以取 \(k(e_1+e_2)\)。商空间本身由 \(U\subseteq V\) 决定；用一个补空间代表它的元素则需要额外选择。上述维数加法公式不要求有限维；\(V\) 有限维时，可以按整数相减得到 \(\dim(V/U)=\dim V-\dim U\)。无限基数的加法不能如此消去，例如 \(\aleph_0=\aleph_0+1=\aleph_0+\aleph_0\)，所以仅凭 \(\dim V\) 和 \(\dim U\) 一般不能恢复商空间的维数。

\subsection{线性映射的核、像与秩}
\label{b2:subsec:0101:03}

\begin{definition}\label{b2:def:linear-map}
设 \(V,W\) 为域 \(k\) 上的向量空间。若映射 \(f:V\to W\) 对任意 \(a,b\in k\) 和 \(u,v\in V\) 满足 \(f(au+bv)=a\cdot f(u)+b\cdot f(v)\)，就称 \(f\) 为\term[xianxing-yingshe]{线性映射}[linear map]。全部这样的映射构成向量空间 \(\operatorname{Hom}_k(V,W)\)，加法和标量乘法逐点定义。若线性映射 \(f\) 是双射，就称它为\term[xianxing-tonggou]{线性同构}[linear isomorphism]；存在这样的映射时记 \(V\cong W\)。
\end{definition}
\symidx{\(\operatorname{Hom}_k(V,W)\)：向量空间之间的线性映射空间}

同构的逆映射也线性：把 \(w_1=f(v_1),w_2=f(v_2)\) 代入线性等式，再施以逆映射即可。线性映射的复合仍线性，恒等映射线性。一个重要的构造办法是，在 \(V\) 的基 \((e_i)\) 上任意指定 \(f(e_i)=w_i\)，再令
\[
f\left(\sum_{i\in F}a_ie_i\right)=\sum_{i\in F}a_iw_i.
\]
基展开的唯一性使定义良好，且给出唯一的线性延拓。这里允许任意指定的是基向量的像；对一般生成组，其线性关系也必须被保持。

\begin{definition}\label{b2:def:kernel-image-rank}
设 \(f:V\to W\) 是 \(k\) 向量空间之间的线性映射。\(f\) 的\term[he]{核}[kernel]、\term[xiang]{像}[image]及\term[zhi]{秩}[rank]分别为
\[
\ker f=\{v\in V\mid f(v)=0\},\qquad
\operatorname{im}f=\{f(v)\mid v\in V\},\qquad
\operatorname{rank}f=\dim_k\operatorname{im}f.
\]
\end{definition}
\symidx{\(\ker f,\operatorname{im}f\)：线性映射的核与像}
\symidx{\(\operatorname{rank}f\)：线性映射的秩}

核与像分别是 \(V,W\) 的子空间：例如 \(f(au+bv)=0\) 保证核的封闭性，\(a\cdot f(u)+b\cdot f(v)=f(au+bv)\) 保证像的封闭性。由 \(f(v)=f(w)\iff v-w\in\ker f\)，\(f\) 单射当且仅当 \(\ker f=0\)。

\begin{proposition}[商空间的映射性质]\label{b2:prop:quotient-linear-universal}
设 \(U\le V\) 是域 \(k\) 上的向量空间 \(V\) 的子空间，\(f:V\to W\) 是 \(k\) 线性映射，\(q:V\to V/U\) 是商映射。存在满足 \(\bar f q=f\) 的线性映射 \(\bar f:V/U\to W\)，当且仅当 \(U\subseteq\ker f\)。存在时它唯一，且 \(\bar f(v+U)=f(v)\)。

条件 \(U\subseteq\ker f\) 保证下图存在唯一的虚线箭头，且两条从 \(V\) 到 \(W\) 的路径相同：
\[
\begin{tikzcd}[column sep=3.5em,row sep=2.5em]
V\arrow[r,"f"]\arrow[d,"q"']&W\\
V/U\arrow[ur,dashed,"\bar f"']&
\end{tikzcd}
\]
\end{proposition}
\begin{proof}
若 \(\bar f q=f\)，对 \(u\in U\) 有 \(f(u)=\bar f(0)=0\)。反过来，若 \(U\subseteq\ker f\)，则 \(v-v'\in U\) 蕴含 \(f(v)=f(v')\)，所给公式与代表元无关。线性性由商空间运算及 \(f\) 的线性性得到。每个陪集都在 \(q\) 的像中，所以条件 \(\bar f q=f\) 已决定全部取值。
\end{proof}

\begin{proposition}[线性映射第一同构定理]\label{b2:prop:linear-first-isomorphism}
设 \(f:V\to W\) 是域 \(k\) 上的向量空间之间的线性映射。映射 \(v+\ker f\mapsto f(v)\) 给出 \(k\) 线性同构 \(V/\ker f\cong\operatorname{im}f\)。若 \(C\le V\) 是 \(\ker f\) 的任意补空间，即 \(V=\ker f\oplus C\)，则限制映射
\[
f|_C:C\longrightarrow\operatorname{im}f
\]
也是 \(k\) 线性同构。
\end{proposition}
\begin{proof}
将\cref{b2:prop:quotient-linear-universal}用于 \(U=\ker f\)，得到以 \(\operatorname{im}f\) 为陪域的线性满射。若 \(f(v)=0\)，则 \(v+\ker f=0\)，所以它的核为零，进而单射。

给定补空间 \(C\)，每个 \(v\in V\) 唯一写为 \(v=n+c\)，其中 \(n\in\ker f,c\in C\)，故 \(f(v)=f(c)\)，说明 \(f|_C\) 满射到像。若 \(c\in C\) 且 \(f(c)=0\)，则 \(c\in C\cap\ker f=0\)，故限制映射也单射。令 \(q:V\to V/\ker f\) 为商映射，它在 \(C\) 上的限制 \(q|_C:C\to V/\ker f\) 同样是同构：上述分解给出每个商类在 \(C\) 中的代表元，而交为零保证代表元唯一。因而 \(f|_C\) 就是 \(q|_C\) 与商到像的同构的复合；选择 \(C\) 只改变代表元，商到像的同构始终是 \(v+\ker f\mapsto f(v)\)。
\end{proof}

\begin{theorem}[秩与核的维数]\label{b2:thm:rank-nullity}
设 \(f:V\to W\) 是域 \(k\) 上的向量空间之间的线性映射，则
\[
\dim_kV=\dim_k\ker f+\operatorname{rank}f.
\]
等式中的加法是基数加法；若 \(V\) 有限维，它就是非负整数的加法。
\end{theorem}
\begin{proof}
由\cref{b2:prop:quotient-basis}，\(\dim_kV=\dim_k\ker f+\dim_k(V/\ker f)\)。由\cref{b2:prop:linear-first-isomorphism}，商空间与像同构；同构把一组基送到一组基，所以 \(\dim_k(V/\ker f)=\dim_k\operatorname{im}f=\operatorname{rank}f\)，得到所述等式。
\end{proof}

这个等式通常称为\term[zhi-lingdu-dingli]{秩--零度定理}[rank--nullity theorem]，其中零度指核的维数。

\ProseParagraphBreak
有限维向量空间的线性自映射 \(f:V\to V\) 因而单射当且仅当满射：核维数为零恰好等价于像维数等于 \(\dim V\)，而有限维空间的同维子空间只能是全空间。在 \(k[t]\) 上，这个结论会失败。乘以 \(t\) 的线性映射单射而不满，其像 \(tk[t]\) 虽是真子空间，维数却仍为 \(\aleph_0\)；由 \(1\mapsto0,t^i\mapsto t^{i-1}\ (i\ge1)\) 逐基向量定义的线性映射满射而不单，核为 \(k\cdot1\)。后一个算子仍满足秩--零度等式 \(\aleph_0=1+\aleph_0\)，却不能消去两边的 \(\aleph_0\) 而断言核为零。

\ProseParagraphBreak
作为有限计算，令 \(f:k^3\to k^2\) 为 \(f(x,y,z)=(x+y,y+z)\)。核由 \((-1,1,-1)\) 生成，而 \(f(e_1)=(1,0),f(e_3)=(0,1)\)，所以秩为二。取 \(C=ke_1+ke_3\)。核向量的第二坐标若为零，它就只能是零，故 \(C\cap\ker f=0\)；再由分解
\[
(x,y,z)=y(-1,1,-1)+(x+y,0,z+y)
\]
可知 \(C\) 是核的一个补空间。由\cref{b2:prop:linear-first-isomorphism}，\(f|_C\) 必同构到像 \(k^2\)；在这个选择下，它就是 \((a,0,b)\mapsto(a,b)\)。这既给出 \(3=1+2\)，也展示了补空间如何为 \(k^3/\ker f\) 的每个商类选出唯一代表元。

% !TeX root = ../../Book2.tex
% 纲要标题：### 1.2 对偶与配对
% 纲要位置：计划纲要.md，第 837 行。

\section{对偶与配对}
\label{b2:sec:0102}

% 纲要内容：
% * 代数对偶及双对偶
% * 对偶映射与有限维自然同构
% * 无限维情形不能直接照搬的结论

一个向量可以用坐标描述，也可以用一族线性函数的取值来辨认。例如，在 \(k^n\) 中分别读取各个坐标，就足以确定向量；一个齐次线性方程组的解空间，则由若干线性函数的共同零点给出。对偶空间把这些线性函数本身看成向量，使关于向量和方程的两种描述可以相互比较。本节的基域始终为 \(k\)，不在向量空间上附加拓扑，因而这里的对偶包含全部线性函数。

\subsection{代数对偶与双对偶}
\label{b2:subsec:0102:01}

\begin{definition}\label{b2:def:algebraic-dual}
设 \(V\) 是域 \(k\) 上的向量空间。若 \(\lambda:V\to k\) 是线性映射，就称 \(\lambda\) 为 \(V\) 上的\term[xianxing-fanhan]{线性泛函}[linear functional]。全体线性泛函组成的空间
\[
V^*\coloneqq\operatorname{Hom}_k(V,k)
\]
称为 \(V\) 的\term[daishu-duiou]{代数对偶}[algebraic dual]，通常简称对偶空间。其中的加法与数乘按取值定义：
\[
(\lambda+\mu)(v)=\lambda(v)+\mu(v),\qquad
(a\lambda)(v)=a\cdot\lambda(v).
\]
再对 \(V^*\) 取对偶，得到\term[shuang-duiou]{双对偶}[double dual]
\(V^{**}\coloneqq(V^*)^*\)。
\end{definition}
\symidx{\(V^*\)：向量空间的代数对偶}
\symidx{\(V^{**}\)：向量空间的双对偶}

这里的运算确实仍给出线性泛函。例如，对 \(a,b\in k\) 及 \(u,v\in V\)，
\[
(\lambda+\mu)(au+bv)
=a(\lambda+\mu)(u)+b(\lambda+\mu)(v).
\]
向量空间公理由 \(k\) 中的相应等式逐点成立。\(V^*\) 的元素是作用在向量上的函数，而 \(V^{**}\) 的元素是作用在这些函数上的线性函数；两层定义中的自变量不同。

\ProseParagraphBreak
在 \(V=k^n\) 中，记 \(\varepsilon_i(x_1,\ldots,x_n)=x_i\)。每个线性泛函都具有唯一形式
\[
\lambda(x_1,\ldots,x_n)=a_1x_1+\cdots+a_nx_n,
\qquad a_i=\lambda(e_i),
\]
其中 \(e_i\) 是第 \(i\) 个标准基向量。这就是把齐次线性方程的左边看成一个向量：系数相加，对应泛函相加。

\begin{proposition}[对偶基]\label{b2:prop:dual-basis}
设 \(V\) 是域 \(k\) 上的向量空间，且有有限基 \(e_1,\ldots,e_n\)。存在唯一的泛函
\(\varepsilon_1,\ldots,\varepsilon_n\in V^*\)，满足
\[
\varepsilon_i(e_j)=\delta_{ij},
\qquad
\delta_{ij}=\begin{cases}1,&i=j,\\0,&i\ne j.\end{cases}
\]
它们构成 \(V^*\) 的一组基。特别地，\(\dim_kV^*=\dim_kV=n\)。
\end{proposition}

\begin{proof}
定义 \(\varepsilon_i(\sum_jx_je_j)=x_i\)。基展开的唯一性保证良定义，坐标对加法与数乘的相容性保证线性。一个线性映射由基向量的像唯一确定，因此这些泛函唯一。对任意 \(\lambda\in V^*\)，在每个基向量上取值可得
\[
\lambda=\sum_{i=1}^n\lambda(e_i)\varepsilon_i.
\]
若 \(\sum_i a_i\varepsilon_i=0\)，在 \(e_j\) 上取值得 \(a_j=0\)。所以它们既生成 \(V^*\)，又线性无关。
\end{proof}

这组 \(\varepsilon_i\) 称为给定基的\term[duiou-ji]{对偶基}[dual basis]。它依赖原来的基；对偶空间本身的定义则没有选择基。例如在 \(k[t]\) 上，取常数项、取 \(t^m\) 的系数、在某个 \(a\in k\) 处取值，都是线性泛函。最后一种泛函满足
\[
\lambda_a\left(\sum_{j=0}^n c_jt^j\right)=\sum_{j=0}^n c_ja^j,
\]
因而与单个系数泛函有不同的含义。这里每个多项式只含有限项，无须赋予无限和任何意义。

\begin{lemma}[泛函的延拓与分离]\label{b2:lem:functional-extension}
设 \(V\) 是域 \(k\) 上的向量空间，\(U\le V\)。
\begin{enumerate}
\item 每个线性泛函 \(\mu:U\to k\) 都能延拓为 \(V\) 上的线性泛函。
\item 若 \(v\notin U\)，则存在 \(\lambda\in V^*\)，使 \(\lambda(u)=0\) 对一切 \(u\in U\) 成立，而 \(\lambda(v)=1\)。
\end{enumerate}
\end{lemma}

\begin{proof}
由\cref{b2:thm:vector-basis-extension}，取 \(U\) 的一组基 \(B\)，再把 \(B\) 扩充为 \(V\) 的基 \(B\cup C\)。在 \(B\) 上保留 \(\mu\) 的取值，在 \(C\) 上指定值为零。每个向量的基展开只有有限项且唯一，故按线性组合取值便定义了所需延拓。

若 \(v\notin U\)，则 \(B\cup\{v\}\) 线性无关：关系 \(av+u=0\) 中若 \(a\ne0\)，就会有 \(v=-a^{-1}u\in U\)。将这组无关向量扩充为 \(V\) 的基，指定 \(v\) 的值为一，其余基向量的值为零，就得到第二个结论。
\end{proof}

延拓一般不唯一：只要 \(U\ne V\)，就可以改变补充基向量上的取值。分离结论则说明，线性泛函足以检测一个向量是否属于指定的子空间；这个事实不限于有限维，只是一般情形的证明使用了基扩充。

\begin{definition}\label{b2:def:linear-pairing}
设 \(V,W\) 是域 \(k\) 上的向量空间。若映射 \(b:V\times W\to k\) 对每个变量分别线性，即对任意 \(a,c\in k\)、\(v,v_1,v_2\in V\) 和 \(w,w_1,w_2\in W\) 满足
\[
\begin{aligned}
b(av_1+cv_2,w)&=a\cdot b(v_1,w)+c\cdot b(v_2,w),\\
b(v,aw_1+cw_2)&=a\cdot b(v,w_1)+c\cdot b(v,w_2),
\end{aligned}
\]
就称 \(b\) 为一个\term[peidui]{配对}[pairing]，也称双线性配对。若只有 \(v=0\) 才满足 \(b(v,w)=0\) 对所有 \(w\in W\) 成立，就称 \(b\) 左非退化；若只有 \(w=0\) 才满足 \(b(v,w)=0\) 对所有 \(v\in V\) 成立，就称 \(b\) 右非退化。两者都成立时称 \(b\) 为\term[fei-tuihua-peidui]{非退化配对}[nondegenerate pairing]。
\end{definition}

例如 \(k^n\times k^n\) 上的 \(b(x,y)=\sum_i x_iy_i\) 非退化，因为用标准基向量作第二个变量，可以逐一读取 \(x\) 的坐标，交换两个变量也是如此。这在任意特征下都成立，并不要求 \(b(x,x)\ne0\) 对每个非零 \(x\) 成立。相反，在 \(k^2\) 上只取 \(b(x,y)=x_1y_1\)，就会漏掉第二个坐标，两侧都退化。\(k^n\) 上的配对也不包含复共轭；当 \(k=\mathbb C\) 时，应将双线性与通常内积中的共轭线性区分开。

\ProseParagraphBreak
每个配对都给出两个线性映射
\[
V\longrightarrow W^*,\quad v\longmapsto b(v,-),
\qquad
W\longrightarrow V^*,\quad w\longmapsto b(-,w).
\]
其中 \(b(v,-)\) 表示把第一个变量固定为 \(v\) 后得到的泛函。左非退化恰好说明第一个映射单射，右非退化恰好说明第二个映射单射。因此，非退化性说的是能否通过配对区分两侧的向量。

\begin{proposition}[评价配对]\label{b2:prop:evaluation-injective}
对任意 \(k\) 向量空间 \(V\)，配对
\[
V\times V^*\longrightarrow k,\qquad(v,\lambda)\longmapsto\lambda(v)
\]
非退化。它给出的\term[pingjia-yingshe]{评价映射}[evaluation map]
\[
j_V:V\longrightarrow V^{**},\qquad j_V(v)(\lambda)=\lambda(v)
\]
是单射线性映射。
\end{proposition}
\symidx{\(j_V\)：从向量空间到双对偶的评价映射}

\begin{proof}
固定 \(v\) 时，\((a\lambda+c\mu)(v)=a\cdot\lambda(v)+c\cdot\mu(v)\)，所以 \(j_V(v)\) 确实属于 \(V^{**}\)。固定 \(\lambda\) 时，由 \(\lambda\) 的线性得到 \(j_V(av+cw)=a\cdot j_V(v)+c\cdot j_V(w)\)。若 \(v\ne0\)，在\cref{b2:lem:functional-extension}中取 \(U=0\)，可得 \(\lambda(v)=1\)，所以 \(j_V(v)\ne0\)，并且配对左非退化。一个非零泛函按定义必在某个向量上非零，因此配对也右非退化。
\end{proof}

\begin{definition}\label{b2:def:vector-annihilator}
设 \(V\) 是域 \(k\) 上的向量空间，\(U\subseteq V\)，\(E\subseteq V^*\)。定义它们关于评价配对 \((v,\lambda)\mapsto\lambda(v)\) 的\term[linghuazi]{零化子}[annihilator]（也称湮灭子）为
\[
\begin{aligned}
U^\circ&\coloneqq\{\lambda\in V^*\mid \lambda(u)=0\text{ 对所有 }u\in U\},\\
{}^\circ E&\coloneqq\{v\in V\mid \lambda(v)=0\text{ 对所有 }\lambda\in E\}.
\end{aligned}
\]
前者位于 \(V^*\)，后者位于 \(V\)，二者都是子空间。本节用圆圈所在的位置区分这两种零化子。
\end{definition}
\termidx[yanmiezi]{湮灭子}
\symidx{\(U^\circ\)：向量子集在对偶空间中的零化子}
\symidx{\({}^\circ E\)：泛函族在原空间中的零化子}

子空间性来自共同零点对线性组合封闭。例如若 \(v,w\in{}^\circ E\)，则每个 \(\lambda\in E\) 都满足 \(\lambda(av+bw)=0\)。在 \(k^3\) 中取 \(U=k e_1\)，就有
\[
U^\circ=k\varepsilon_2+k\varepsilon_3,
\qquad
{}^\circ(k\varepsilon_2+k\varepsilon_3)=k e_1.
\]
向量生成的空间越大，在它上面恒为零的泛函越少。因此 \(U_1\subseteq U_2\) 蕴含 \(U_2^\circ\subseteq U_1^\circ\)；另一侧同样反转包含关系。

\subsection{对偶映射与有限维自然同构}
\label{b2:subsec:0102:02}

给定线性映射 \(f:V\to W\)，一个 \(W\) 上的线性方程可以拉回为 \(V\) 上的线性方程：把原方程中的变量换成 \(f(v)\) 即可。这个简单的代入改变了映射方向。

\begin{definition}\label{b2:def:dual-map}
设 \(f:V\to W\) 是域 \(k\) 上的向量空间之间的线性映射。\(f\) 的\term[duiou-yingshe]{对偶映射}[dual map]定义为
\[
f^*:W^*\longrightarrow V^*,\qquad f^*(\lambda)=\lambda\circ f.
\]
\end{definition}
\symidx{\(f^*\)：线性映射的对偶映射}

因为线性映射的复合仍线性，\(f^*(\lambda)\) 属于 \(V^*\)；又因为
\((a\lambda+b\mu)\circ f=a(\lambda\circ f)+b(\mu\circ f)\)，\(f^*\) 本身也是线性映射。注意其定义域为 \(W^*\)，方向与 \(f\) 相反。

\begin{proposition}[对偶映射的复合、核与像]\label{b2:prop:dual-map}
设 \(V,W,Z\) 是域 \(k\) 上的向量空间，\(f:V\to W\)、\(g:W\to Z\) 是线性映射，则
\[
(g\circ f)^*=f^*\circ g^*,\qquad
(\operatorname{id}_V)^*=\operatorname{id}_{V^*},
\]
并且
\begin{equation}\label{b2:eq:dual-kernel-image}
\ker f^*=(\operatorname{im}f)^\circ,
\qquad
\operatorname{im}f^*=(\ker f)^\circ.
\end{equation}
特别地，\(f\) 单射当且仅当 \(f^*\) 满射；\(f\) 满射当且仅当 \(f^*\) 单射。以上结论不要求有限维。
\end{proposition}

\begin{proof}
对 \(\lambda\in Z^*\)，直接计算
\[
(g\circ f)^*(\lambda)=\lambda\circ g\circ f
=f^*\bigl(g^*(\lambda)\bigr).
\]
恒等映射的结论也是逐个泛函代入定义即得。

\(\lambda\in W^*\) 属于 \(\ker f^*\)，恰好表示 \(\lambda\bigl(f(v)\bigr)=0\) 对每个 \(v\in V\) 成立，这就是第一个零化子公式。若 \(\mu=f^*(\lambda)\)，则 \(v\in\ker f\) 时 \(\mu(v)=\lambda(0)=0\)，所以 \(\operatorname{im}f^*\subseteq(\ker f)^\circ\)。反过来，给定在 \(\ker f\) 上为零的 \(\mu\in V^*\)，在 \(\operatorname{im}f\) 上定义
\[
\lambda_0\bigl(f(v)\bigr)=\mu(v).
\]
若 \(f(v)=f(v')\)，则 \(v-v'\in\ker f\)，故 \(\mu(v)=\mu(v')\)，证明了良定义。\(f\) 与 \(\mu\) 的线性给出 \(\lambda_0\) 的线性。由\cref{b2:lem:functional-extension}把 \(\lambda_0\) 延拓到 \(W\) 上的 \(\lambda\)，便有 \(f^*(\lambda)=\mu\)。

若 \(f\) 单射，则 \(\ker f=0\)，上述像公式给出 \(\operatorname{im}f^*=V^*\)。若 \(f^*\) 满射，而 \(0\ne v\in\ker f\)，则 \(V^*\) 中每个泛函都在 \(v\) 处为零，与\cref{b2:lem:functional-extension}的分离结论矛盾。若 \(f\) 满射，则核公式给出 \(\ker f^*=W^\circ=0\)。若 \(\operatorname{im}f\ne W\)，取 \(w\notin\operatorname{im}f\)，同一分离结论给出一个非零泛函，在 \(\operatorname{im}f\) 上为零而在 \(w\) 处为一；它属于 \(\ker f^*\)，所以 \(f^*\) 不单射。
\end{proof}

例如子空间包含映射 \(\iota:U\hookrightarrow V\) 的对偶映射，就是限制映射 \(V^*\to U^*\)。它的满射性是在说泛函能够延拓，其核为 \(U^\circ\)。又如商映射 \(q:V\to V/U\) 的对偶映射是单射，并把 \((V/U)^*\) 识别为 \(U^\circ\)：商空间上的泛函与原空间中在 \(U\) 上为零的泛函一一对应。这里 \(q^*(\lambda)(v)=\lambda(v+U)\)，所以这个识别也不需要选择补空间。

\ProseParagraphBreak
对偶映射在矩阵上表现为转置。设 \(V,W\) 的基分别为 \(e_1,\ldots,e_n\) 与 \(w_1,\ldots,w_m\)，对偶基分别为 \(\varepsilon_i\) 与 \(\eta_j\)。按列记 \(f\) 的矩阵 \(A=(a_{ji})\)，即 \(f(e_i)=\sum_j a_{ji}w_j\)。于是
\[
f^*(\eta_j)(e_i)=a_{ji},\qquad
f^*(\eta_j)=\sum_{i=1}^n a_{ji}\varepsilon_i.
\]
故 \(f^*\) 在这两组对偶基下的矩阵为 \(A^{\mathsf T}\)。例如
\[
f:k^3\longrightarrow k^2,\quad f(x,y,z)=(x+z,y+z)
\]
的矩阵及其对偶矩阵分别为
\[
A=\begin{pmatrix}1&0&1\\0&1&1\end{pmatrix},\qquad
A^{\mathsf T}=\begin{pmatrix}1&0\\0&1\\1&1\end{pmatrix}.
\]
\(W\) 上的泛函 \((s,t)\mapsto as+bt\) 拉回后是 \(ax+by+(a+b)z\)，与转置矩阵的计算一致。这一说明也解释了为什么复合取对偶会逆序：矩阵转置满足 \((BA)^{\mathsf T}=A^{\mathsf T}B^{\mathsf T}\)。

\begin{theorem}[有限维双对偶同构]\label{b2:thm:double-dual}
设 \(V,W\) 是域 \(k\) 上的向量空间，\(j_V:V\to V^{**}\) 与 \(j_W:W\to W^{**}\) 是评价映射。若 \(V\) 有限维，则 \(j_V\) 是同构。对任意线性映射 \(f:V\to W\)，不论维数是否有限，总有
\begin{equation}\label{b2:eq:double-dual-natural}
j_W\circ f=f^{**}\circ j_V,
\qquad f^{**}\coloneqq(f^*)^*.
\end{equation}
因此，有限维空间与其双对偶的上述同构既不依赖基，又与所有线性映射相容。
\end{theorem}

\begin{proof}
\cref{b2:prop:evaluation-injective}已经给出 \(j_V\) 的单射性。取 \(V\) 的基 \(e_1,\ldots,e_n\) 及其对偶基 \(\varepsilon_1,\ldots,\varepsilon_n\)。对任意 \(F\in V^{**}\)，令
\[
v=\sum_{i=1}^n F(\varepsilon_i)e_i.
\]
\cref{b2:prop:dual-basis}给出 \(\lambda=\sum_i\lambda(e_i)\varepsilon_i\)，故
\[
F(\lambda)=\sum_{i=1}^n\lambda(e_i)F(\varepsilon_i)=\lambda(v)=j_V(v)(\lambda).
\]
所以每个 \(F\) 都来自某个 \(v\)，证明了满射性。基只用来证明满射；\(j_V\) 的定义本身没有用到这组基。

对一般的 \(f:V\to W\)，给定 \(v\in V\) 与 \(\lambda\in W^*\)，有
\[
\begin{aligned}
\bigl(f^{**}\circ j_V\bigr)(v)(\lambda)
&=j_V(v)\bigl(f^*(\lambda)\bigr)\\
&=(\lambda\circ f)(v)\\
&=\lambda\bigl(f(v)\bigr)\\
&=\bigl(j_W\circ f\bigr)(v)(\lambda).
\end{aligned}
\]
两个 \(W^*\) 上的泛函在每个 \(\lambda\) 处取值相同，即得\eqref{b2:eq:double-dual-natural}。
\end{proof}

这里的“自然”具体指\eqref{b2:eq:double-dual-natural}：可以先把向量经 \(f\) 送到 \(W\)，再取评价泛函；也可以先取评价泛函，再施加 \(f^{**}\)，结果相同。\chapref{b2:ch:05}将给这种相容性统一的名称；眼下只需掌握这个等式。与此不同，给有限维空间选择一组基，再把各基向量送到相应的对偶基，虽然也得到 \(V\cong V^*\)，却使用了额外选择。\(V\) 与 \(V^{**}\) 之间有已经指定的评价映射，不应与后一种依赖基的同构混为一谈。

\begin{corollary}[有限维零化子的维数]\label{b2:cor:annihilator-dimension}
设 \(V\) 是域 \(k\) 上的 \(n\) 维向量空间，\(U\subseteq V\)、\(E\subseteq V^*\) 为子空间，则
\[
\begin{aligned}
\dim U^\circ&=n-\dim U,&{}^\circ(U^\circ)&=U,\\
\dim{}^\circ E&=n-\dim E,&({}^\circ E)^\circ&=E.
\end{aligned}
\]
其中 \({}^\circ(U^\circ)=U\) 对无限维空间也成立。
\end{corollary}

\begin{proof}
设 \(u_1,\ldots,u_r\) 是 \(U\) 的基，把它扩充为 \(V\) 的基 \(u_1,\ldots,u_n\)，记相应对偶基为 \(\varepsilon_1,\ldots,\varepsilon_n\)。泛函在 \(U\) 上为零当且仅当前 \(r\) 个系数为零，所以 \(\varepsilon_{r+1},\ldots,\varepsilon_n\) 是 \(U^\circ\) 的基。对任意维数都有 \(U\subseteq{}^\circ(U^\circ)\)；若 \(v\notin U\)，\cref{b2:lem:functional-extension}给出 \(\lambda\in U^\circ\) 且 \(\lambda(v)=1\)，排除了 \(v\in{}^\circ(U^\circ)\)。

把刚才的维数公式应用到 \(V^*\) 的子空间 \(E\)，可得在 \(V^{**}\) 中对 \(E\) 恒为零的泛函所成空间维数为 \(n-\dim E\)。由\cref{b2:thm:double-dual}，\(j_V\) 将 \({}^\circ E\) 同构到这个空间，因此 \(\dim{}^\circ E=n-\dim E\)。再次使用第一条维数公式，得到 \(\dim({}^\circ E)^\circ=\dim E\)。定义直接给出 \(E\subseteq({}^\circ E)^\circ\)，同维有限维子空间的包含只能是相等。
\end{proof}

于是，有限维时可以在子空间与它所满足的全部齐次线性方程之间来回转换。比如 \(k^3\) 中的平面 \(x+y+z=0\) 是泛函 \(\varepsilon_1+\varepsilon_2+\varepsilon_3\) 的核；它的零化子恰好由这一个泛函生成。若只列出一份方程表，对偶空间关注的是这些方程张成的空间，重复一条方程或对方程作可逆行变换都不改变它。

\ProseParagraphBreak
若 \(V,W\) 有限维，结合\eqref{b2:eq:dual-kernel-image}、\cref{b2:cor:annihilator-dimension}及\cref{b2:thm:rank-nullity}，还有
\[
\operatorname{rank}f^*
=\dim(\ker f)^\circ
=\dim V-\dim\ker f
=\operatorname{rank}f.
\]
因此一个矩阵与其转置的秩相等。这里先描述零化子，再数维数；并没有事先把行秩等于列秩作为依据。

\ProseParagraphBreak
上述结论还说明，两个有限维空间之间的非退化配对会给出 \(V\cong W^*\) 与 \(W\cong V^*\)。事实上，非退化性给出的两个单射蕴含 \(\dim V\le\dim W\) 及 \(\dim W\le\dim V\)，故两侧同维，这两个单射都是同构。无限维时，仅有两个单射不足以使这两个特定映射满射；评价配对很快就会给出反例。

\subsection{无限维空间的对偶}
\label{b2:subsec:0102:03}

有限维时，对偶基能够表达每个泛函；无限维时，基向量仍决定泛函，却不再保证泛函是有限多个坐标泛函的线性组合。区别可以从直和与直积的定义直接看出。

\ProseParagraphBreak
给定一族 \(k\) 向量空间 \((V_i)_{i\in I}\)，\term[zhiji]{直积}[direct product] \(\prod_{i\in I}V_i\) 包含所有向量族 \((v_i)_{i\in I}\)，运算逐坐标进行。\term[zhihe]{直和}[direct sum] \(\bigoplus_{i\in I}V_i\) 是其中满足 \(v_i\ne0\) 的指标只有有限多个的子空间。因此直和中的向量是有限个分量向量之和，直积中的元素则不必有有限支集。记 \(\iota_i:V_i\to\bigoplus_jV_j\) 为放入第 \(i\) 个坐标的线性映射。
\symidx{\(\prod_iV_i\)：向量空间族的直积}
\symidx{\(\bigoplus_iV_i\)：向量空间族的直和}

\begin{proposition}[直和的对偶]\label{b2:prop:dual-direct-sum}
设 \((V_i)_{i\in I}\) 是域 \(k\) 上的向量空间族，\(\iota_i:V_i\to\bigoplus_{j\in I}V_j\) 是典范嵌入。限制到各个分量给出 \(k\) 线性同构
\[
\Phi:\left(\bigoplus_{i\in I}V_i\right)^*
\longrightarrow\prod_{i\in I}V_i^*,
\qquad
\Phi(\lambda)=(\lambda\circ\iota_i)_{i\in I}.
\]
\end{proposition}

\begin{proof}
\(\Phi\) 逐分量线性。给定任意泛函族 \((\lambda_i)\in\prod_iV_i^*\)，定义
\[
\lambda\bigl((v_i)_{i\in I}\bigr)=\sum_{i\in I}\lambda_i(v_i)
\quad\text{对 }(v_i)\in\bigoplus_iV_i.
\]
\((v_i)\) 只有有限多个非零分量，所以这是有限和，且逐项相加与数乘证明 \(\lambda\) 线性。它在第 \(i\) 个分量上的限制为 \(\lambda_i\)，故 \(\Phi\) 满射。每个直和中的向量都是有限个分量向量之和，所以全部限制为零的泛函只能为零，证明了单射性。
\end{proof}

特别地，令 \(\mathbb N=\{1,2,\ldots\}\)，考虑
\[
V=k^{(\mathbb N)}
=\{(x_i)_{i\ge1}\mid x_i\in k,\ x_i\ne0\text{ 仅对有限多个 }i\},
\qquad
k^{\mathbb N}=\prod_{i\ge1}k.
\]
标准基向量 \(e_i\) 构成 \(V\) 的基。上面的同构把 \(V^*\) 识别为 \(k^{\mathbb N}\)：数列 \(a=(a_i)\) 对应泛函
\[
\lambda_a(x)=\sum_{i\ge1}a_ix_i.
\]
允许 \(a\) 有无限多个非零分量，因为被代入的 \(x\) 只有有限支集。不能反过来要求 \(a\) 也有有限支集，否则会漏掉合法的线性泛函。
\symidx{\(k^{(\mathbb N)}\)：有限支集数列空间}
\symidx{\(k^{\mathbb N}\)：全体数列组成的直积空间}

\ProseParagraphBreak
取 \(a_i=1\) 对所有 \(i\) 成立，就得到坐标和泛函
\(\lambda_a(x)=\sum_i x_i\)。它不能是有限多个坐标泛函 \(\varepsilon_i(x)=x_i\) 的线性组合：假如只用了指标集合 \(J\) 中的坐标，取 \(m\notin J\)，便会在 \(e_m\) 上得到零，但 \(\lambda_a(e_m)=1\)。所以 \(\varepsilon_i\) 虽然线性无关，却不是 \(V^*\) 的基；将其称为“无限对偶基”而沿用有限维中的基展开，会造成错误。

\begin{example}[双对偶中不是评价的泛函]\label{b2:ex:infinite-double-dual}
设 \(V=k^{(\mathbb N)}\)。评价映射 \(j_V:V\to V^{**}\) 不是满射。
\end{example}

\begin{proof}
按\cref{b2:prop:dual-direct-sum}把 \(V^*\) 看成 \(k^{\mathbb N}\)。其中有限支集数列组成子空间 \(D=k^{(\mathbb N)}\)，常数列 \(\mathbf1=(1,1,\ldots)\) 不属于 \(D\)。因此 \(D+k\mathbf1\) 中的每个元素有唯一表达 \(d+c\mathbf1\)，可以定义线性泛函
\[
F_0:D+k\mathbf1\longrightarrow k,
\qquad F_0(d+c\mathbf1)=c.
\]
唯一性来自 \(D\cap k\mathbf1=0\)：一个非零常数列不可能有有限支集。由\cref{b2:lem:functional-extension}，将 \(F_0\) 延拓为 \(k^{\mathbb N}\) 上的泛函 \(F\)。于是 \(F\in V^{**}\)，\(F\) 在 \(D\) 上为零，而 \(F(\mathbf1)=1\)。

假若 \(F=j_V(v)\)，其中 \(v=(v_i)\in V\)，则把坐标泛函 \(\varepsilon_i\) 看成 \(k^{\mathbb N}\) 中的第 \(i\) 个标准基向量，它属于 \(D\)，故
\[
v_i=\varepsilon_i(v)=j_V(v)(\varepsilon_i)=F(\varepsilon_i)=0
\]
对所有 \(i\) 成立。因此 \(v=0\)，这又会给出 \(F(\mathbf1)=0\)，与构造矛盾。
\end{proof}

这个反例在任意基域上有效。延拓 \(F\) 的存在使用基扩充，并没有给出一种从任意数列通过有限次坐标运算计算 \(F\) 的公式。我们只需它在 \(D\) 上为零、在 \(\mathbf1\) 处为一，就足以证明它不可能是某个有限支集向量的评价。非退化评价配对仍然存在，但它给出的映射 \(V\to V^{**}\) 不满，故有限维非退化配对所带来的同构不能照搬到这里。

\ProseParagraphBreak
零化子的两侧也因此出现差别。在这个例子中，令 \(E\subseteq V^*\) 为所有坐标泛函的线性生成空间，即上面的 \(D\)。坐标能够分辨 \(V\) 中的向量，所以 \({}^\circ E=0\)，进而 \(({}^\circ E)^\circ=V^*\)；然而 \(E\ne V^*\)。相比之下，\cref{b2:cor:annihilator-dimension}中由泛函分离证明的 \({}^\circ(U^\circ)=U\) 对 \(V\) 的任意子空间仍然成立。区分这两个等式，比把所有对偶结论一概限制在有限维更准确。

% !TeX root = ../../Book2.tex
% 纲要标题：### 1.3 线性算子
% 纲要位置：计划纲要.md，第 843 行。

\section{线性算子}
\label{b2:sec:0103}

给定一个向量及一个线性算子，可以反复施加算子，得到 \(v,Tv,T^2v,\ldots\)。这些向量之间的关系记录了算子的一部分结构。本节把关系写成多项式，说明一个算子如何使整个多项式环作用在空间上。这将是下一章从向量空间转向模的一个主要例子。

% 纲要内容（仅作写作计划，尚非教材正文）：
%
% * 特征子空间、广义特征子空间及互素核分解
% * 最小多项式与循环向量
% * 将线性算子看成多项式环的作用；限制与商算子的两层零化关系
%

\subsection{特征子空间与广义特征子空间}
\label{b2:subsec:0103:01}

\begin{definition}\label{b2:def:linear-operator}
设 \(V\) 是域 \(k\) 上的向量空间。给定线性映射 \(T:V\to V\)，称 \(T\) 为 \(V\) 上的\term[xianxing-suanzi]{线性算子}[linear operator]，也称线性自同态。记
\(\operatorname{End}_k(V)=\operatorname{Hom}_k(V,V)\)。其中加法逐点定义，乘法采用复合，约定 \(ST=S\circ T\)，即先施加右侧的 \(T\)。恒等算子记为 \(I_V\)，并令 \(T^0=I_V\)。
\end{definition}

\symidx{\(\operatorname{End}_k(V)\)：向量空间 \(V\) 的线性自同态环}

复合满足结合律，又对逐点加法满足左右分配律，所以 \(\operatorname{End}_k(V)\) 是含幺环。它通常不交换；标量算子 \(aI_V\) 则与每个线性算子交换，因为 \(T(av)=a\cdot T(v)\)。以下除单独讨论外，\(V\) 有限维且非零，记 \(n=\dim_kV\)。

\begin{definition}\label{b2:def:eigen-generalized}
设 \(T\) 是域 \(k\) 上的向量空间 \(V\) 的线性算子，\(\lambda\in k\)。若 \(0\ne v\in V\) 且 \(Tv=\lambda v\)，就称 \(v\) 为 \(T\) 对应于 \(\lambda\) 的\term[tezheng-xiangliang]{特征向量}[eigenvector]，并称 \(\lambda\) 为 \(T\) 的\term[tezheng-zhi]{特征值}[eigenvalue]。对应于 \(\lambda\) 的\term[tezheng-zi-kongjian]{特征子空间}[eigenspace]和\term[guangyi-tezheng-zi-kongjian]{广义特征子空间}[generalized eigenspace]为
\[
E_\lambda(T)=\ker(T-\lambda I_V),\qquad
G_\lambda(T)=\bigcup_{r\ge1}\ker\bigl((T-\lambda I_V)^r\bigr).
\]
\end{definition}
\symidx{\(E_\lambda(T),G_\lambda(T)\)：特征子空间与广义特征子空间}

这些核依次包含，故其并是子空间。每个核都被 \(T\) 保持：\(T\) 与 \((T-\lambda I_V)^r\) 交换。一般地，设 \(T\) 是 \(k\) 向量空间 \(V\) 上的算子。若子空间 \(U\le V\) 满足 \(T(U)\subseteq U\)，就称 \(U\) 为 \(T\) 的\term[bubian-zi-kongjian]{不变子空间}[invariant subspace]。

\begin{proposition}\label{b2:prop:kernel-power-stability}
设 \(V\) 是域 \(k\) 上的 \(n\) 维向量空间，\(N\in\operatorname{End}_k(V)\)，则核链 \(\ker(N^r)\) 至迟在 \(r=n\) 时稳定，即 \(\ker(N^n)=\ker(N^{n+1})=\cdots\)。
\end{proposition}
\begin{proof}
令 \(K_r=\ker(N^r)\)，其中 \(K_0=0\)。若 \(K_r=K_{r+1}\)，对 \(v\in K_{r+2}\) 有 \(Nv\in K_{r+1}=K_r\)，所以 \(v\in K_{r+1}\)。因此一旦相邻两项相等，以后各项都相等。每一次严格包含至少使维数增加一；从维数零开始至多增加 \(n\) 次，故 \(K_n=K_{n+1}=\cdots\)。
\end{proof}

将命题中的 \(N\) 取为 \(T-\lambda I_V\)，就得到
\[
\begin{gathered}
\ker(T-\lambda I_V)\subseteq\ker\bigl((T-\lambda I_V)^2\bigr)\subseteq\cdots,\\
G_\lambda(T)=\ker\bigl((T-\lambda I_V)^n\bigr).
\end{gathered}
\]
因此在有限维空间中，不必取无穷多个核的并，只需计算核链稳定后的一项。

\ProseParagraphBreak
例如在 \(k^2\) 的基 \(e_1,e_2\) 下取
\[
Te_1=\lambda e_1,\qquad Te_2=e_1+\lambda e_2.
\]
此时 \(N=T-\lambda I_V\) 满足 \(Ne_2=e_1,Ne_1=0\)。所以 \(E_\lambda=ke_1\)，而 \(G_\lambda=V\)。对应于 \(\lambda\) 的广义特征向量，记录反复施加 \(T-\lambda I_V\) 后何时变为零；普通特征向量在第一次施加该算子时便变为零。

\ProseParagraphBreak
不同特征值的特征向量线性无关。确切地说，若存在非平凡关系，取其中非零项数最少的一条，写成
\[
\sum_{j=1}^ra_jv_j=0,
\]
其中各 \(\lambda_j\) 不同且 \(a_j\ne0\)，则 \(r\ge2\)。施加 \(T-\lambda_rI_V\)，得到少一项的关系，其余系数 \(a_j(\lambda_j-\lambda_r)\) 仍非零，矛盾。不过特征向量未必生成全空间，上面的二维例子已经说明这一点；基域中还可能根本没有特征值。

\subsection{最小多项式与循环向量}
\label{b2:subsec:0103:02}

对于 \(p(t)=\sum_{i=0}^da_it^i\in k[t]\)，定义
\[
p(T)=\sum_{i=0}^da_iT^i.
\]
标量算子与 \(T\) 交换，因而直接展开有限和可得
\[
(p+q)(T)=p(T)+q(T),\qquad (pq)(T)=p(T)q(T),\qquad 1(T)=I_V.
\]
这些公式说明求值映射
\[
\Phi_T:k[t]\longrightarrow\operatorname{End}_k(V),\qquad p\longmapsto p(T)
\]
保持加法、乘法和单位。记它的像为
\[
k[T]\coloneqq\operatorname{im}\Phi_T=\{\,p(T)\mid p\in k[t]\,\}.
\]
\symidx{\(k[T]\)：由线性算子生成的多项式算子代数}
这里 \(k[t]\) 中的元素是形式多项式，\(k[T]\) 中的元素则是实际作用于 \(V\) 的算子；不同多项式可能给出同一个算子，例如 \(T^2=0\) 时 \(t^2\) 与零多项式的像相同。虽然 \(\operatorname{End}_k(V)\) 通常不交换，\(k[T]\) 中的算子却彼此交换。

\begin{definition}\label{b2:def:minimal-polynomial}
设 \(V\) 是域 \(k\) 上的向量空间，\(T\in\operatorname{End}_k(V)\)，\(p\in k[t]\)。若 \(p(T)=0\)，就称 \(p\) \term[linghua]{零化}[annihilate]算子 \(T\)。如果存在非零零化多项式，就把其中次数最小的首一多项式称为 \(T\) 的\term[zuixiao-duoxiangshi]{最小多项式}[minimal polynomial]，记作 \(\mu_T(t)\)。
\end{definition}
\symidx{\(\mu_T\)：线性算子的最小多项式}

零向量空间上的恒等算子也是零算子，因此这时 \(\mu_T=1\)。

\begin{proposition}\label{b2:prop:minimal-polynomial-divisibility}
设 \(V\) 是域 \(k\) 上的有限维向量空间，\(T\in\operatorname{End}_k(V)\)，则 \(T\) 有唯一的最小多项式 \(\mu_T\)，且对任意 \(p\in k[t]\) 有
\[
p(T)=0\quad\Longleftrightarrow\quad \mu_T\mid p.
\]
\end{proposition}
\begin{proof}
若 \(V=0\)，任意 \(p(T)\) 都是零算子，而 \(\mu_T=1\)，结论立即成立。以下设 \(V\ne0\)，令 \(n=\dim_kV\)。在一组基上指定像，便以 \(n^2\) 个标量确定一个自同态，所以 \(\operatorname{End}_k(V)\) 维数为 \(n^2\)。因此 \(I_V,T,\ldots,T^{n^2}\) 线性相关，存在非零零化多项式。选次数最小者，再除以其首项系数，得到首一的 \(\mu_T\)。它的次数至少为一，因为非零常数算子不是零算子。

对任意 \(p\) 作带余除法 \(p=qm_T+r\)，其中 \(r=0\) 或 \(\deg r<\deg \mu_T\)。若 \(p(T)=0\)，则 \(r(T)=0\)；最小性迫使 \(r=0\)。反向蕴含由乘法公式立即得到。另一同次数的首一最小零化多项式必须为 \(\mu_T\) 的常数倍，首一条件又使这个倍数为一。
\end{proof}

上述证明中的 \(n^2\) 只给出存在非零零化多项式的粗界。
\cref{b2:prop:operator-characteristic-minimal}将在第四章证明 \(\mu_T\mid\chi_T\)，从而得到更强的 \(\deg \mu_T\le n\)。这里也已确定求值映射的核：\(\ker\Phi_T=(\mu_T)\)，即 \(p(T)=q(T)\) 当且仅当 \(\mu_T\mid(p-q)\)。

\begin{corollary}\label{b2:cor:eigenvalue-minimal-root}
设 \(V\ne0\) 是域 \(k\) 上的有限维向量空间，\(T\in\operatorname{End}_k(V)\)，\(\lambda\in k\)。则 \(\lambda\) 是 \(T\) 的特征值，当且仅当 \(\mu_T(\lambda)=0\)。
\end{corollary}
\begin{proof}
若 \(0\ne v\in V\) 满足 \(Tv=\lambda v\)，则 \(0=\mu_T(T)v=\mu_T(\lambda)v\)，故 \(\mu_T(\lambda)=0\)。反过来，若 \(\mu_T=(t-\lambda)q\) 而 \(T-\lambda I_V\) 单射，由 \((T-\lambda I_V)q(T)=0\) 可得 \(q(T)=0\)，与 \(\mu_T\) 的最小性矛盾。因此 \(T-\lambda I_V\) 有非零核。
\end{proof}

\begin{definition}\label{b2:def:vector-minimal-polynomial}
设 \(V\) 是域 \(k\) 上的向量空间，\(T\in\operatorname{End}_k(V)\)，\(v\in V\)。若存在非零多项式 \(p\in k[t]\) 使 \(p(T)v=0\)，则在这样的多项式中取次数最小的首一者，称为 \(v\) 相对于 \(T\) 的最小多项式，记作 \(\mu_{T,v}\)。特别地，\(\mu_{T,0}=1\)。
\end{definition}
\symidx{\(\mu_{T,v}\)：向量相对于算子的最小多项式}

带余除法表明 \(\mu_{T,v}\) 唯一，且 \(p(T)v=0\) 当且仅当 \(\mu_{T,v}\mid p\)。若 \(V\) 有限维，对非零 \(v\) 可利用 \(v,Tv,\ldots,T^{\dim_kV}v\) 的线性相关性得到非零零化多项式，因此每个向量都有 \(\mu_{T,v}\)。只要 \(\mu_T\) 存在，\(\mu_T(T)v=0\) 还给出 \(\mu_{T,v}\mid \mu_T\)。

\begin{definition}\label{b2:def:cyclic-vector}
设 \(V\) 是域 \(k\) 上的向量空间，\(T\in\operatorname{End}_k(V)\)，\(v\in V\)。由 \(v\) 产生的\term[xunhuan-zi-kongjian]{循环子空间}[cyclic subspace]是
\[
k[T]v=\{p(T)v\mid p\in k[t]\}.
\]
它是所有 \(T^jv\ (j\ge0)\) 的线性生成空间，也是 \(k[T]\) 中的算子作用于 \(v\) 所得到的全部向量组成的空间。若 \(k[T]v=V\)，就称 \(v\) 为 \(T\) 的\term[xunhuan-xiangliang]{循环向量}[cyclic vector]。
\end{definition}
\symidx{\(k[T]v\)：由向量生成的循环子空间}

映射 \(k[t]\to k[T]v\)、\(p\mapsto p(T)v\) 的核记录了 \(v\) 满足的全部多项式关系。有非零关系时，这个核是 \((\mu_{T,v})\)；没有非零关系时，核为零。与 \(\Phi_T\) 的核不同，这里只要求零化 \(v\)，不要求零化整个空间。

\begin{proposition}\label{b2:prop:cyclic-basis}
设 \(V\ne0\) 是域 \(k\) 上的有限维向量空间，\(T\in\operatorname{End}_k(V)\)，\(0\ne v\in V\)。令 \(\mu_{T,v}\) 为零化 \(v\) 的首一最小多项式，\(d=\deg \mu_{T,v}\)。此时 \(v,Tv,\ldots,T^{d-1}v\) 是 \(k[T]v\) 的基。若 \(v\) 是循环向量，则 \(\mu_T=\mu_{T,v}\)，从而 \(\deg \mu_T=\dim_k V\)。
\end{proposition}
\begin{proof}
这些向量若相关，就有一个次数小于 \(d\) 的非零多项式零化 \(v\)，矛盾。对任意 \(p\)，用 \(\mu_{T,v}\) 除去高次项，得到 \(p(T)v=r(T)v\)、\(\deg r<d\)，所以它们生成循环子空间。

若这个空间为 \(V\)，则每个向量写成 \(p(T)v\)。算子多项式彼此交换给出
\[
\mu_{T,v}(T)p(T)v=p(T)\mu_{T,v}(T)v=0,
\]
故 \(\mu_{T,v}\) 零化整个算子，上一命题给出 \(\mu_T\mid \mu_{T,v}\)。另一方面，\(\mu_T(T)v=0\) 给出 \(\mu_{T,v}\mid \mu_T\)。两者均首一，故相等。
\end{proof}

循环向量不总存在。\(n>1\) 时，标量算子 \(T=\lambda I_V\) 的每个非零循环子空间都只有一维。相反，前面的二维例子以 \(e_2\) 为循环向量，因为 \(e_2,Te_2\) 生成 \(e_1,e_2\)；其最小多项式为 \((t-\lambda)^2\)。

\begin{proposition}[互素多项式给出的分解]\label{b2:prop:coprime-operator-decomposition}
设 \(V\) 是域 \(k\) 上的向量空间，\(T\in\operatorname{End}_k(V)\)。若 \(a,b\in k[t]\) 互素且 \(a(T)b(T)=0\)，则
\[
V=\ker a(T)\oplus\ker b(T).
\]
\end{proposition}
\begin{proof}
由第一卷第11.2节的 Euclid 算法命题\footnote{欧几里得（希腊文 \OriginalName{Εὐκλείδης}，约公元前325-约公元前265），古希腊数学家，著有《几何原本》，其中系统讨论了辗转相除法。}给出的 Bézout 等式\footnote{贝祖（Étienne Bézout，1730-1783），法国数学家，研究代数方程的消元与结式，并著有数学教材。}，互素多项式 \(a,b\) 的最大公因子为单位，故可取 \(r,s\in k[t]\)，使 \(ra+sb=1\)。于是
\[
v=r(T)a(T)v+s(T)b(T)v.
\]
第一项被 \(b(T)\) 零化，第二项被 \(a(T)\) 零化，故两核之和为全空间。若 \(v\) 同时在两核中，同一等式给出 \(v=0\)，所以和是直和。核由算子多项式定义，也都是 \(T\) 不变子空间。
\end{proof}

例如 \((T-\lambda I_V)(T-\mu I_V)=0\) 且 \(\lambda\ne\mu\) 时，每个向量唯一分成两个特征子空间中的分量。最小多项式中出现重复的一次因子时，同一互素分解把普通特征子空间换成广义特征子空间。

\begin{corollary}\label{b2:cor:generalized-eigenspace-decomposition}
设 \(V\ne0\) 是域 \(k\) 上的有限维向量空间，\(T\in\operatorname{End}_k(V)\)。若
\[
\mu_T(t)=\prod_{i=1}^s(t-\lambda_i)^{a_i},
\qquad \lambda_i\in k\;\text{两两不同},\quad a_i\ge1,
\]
则
\[
V=\bigoplus_{i=1}^sG_{\lambda_i}(T),\qquad
G_{\lambda_i}(T)=\ker\bigl((T-\lambda_iI_V)^{a_i}\bigr).
\]
各个直和分量都是 \(T\) 不变子空间。
\end{corollary}
\begin{proof}
令 \(q_i=(t-\lambda_i)^{a_i}\)。这些多项式两两互素，且 \(\prod_iq_i(T)=\mu_T(T)=0\)。先证明一般的直和公式：只要 \(q_1,\ldots,q_s\) 两两互素且乘积零化算子，就有 \(V=\bigoplus_i\ker q_i(T)\)。对因子个数 \(s\) 归纳：\(s=1\) 时 \(q_1(T)=0\)；\(s>1\) 时先将 \(q_1\) 与 \(q_2\cdots q_s\) 代入\cref{b2:prop:coprime-operator-decomposition}，得到
\[
V=\ker q_1(T)\oplus\ker\bigl(q_2(T)\cdots q_s(T)\bigr).
\]
第二个分量被 \(T\) 保持，在其上应用归纳假设便得到其余各核的直和。这里每个 \(\ker q_i(T)\ (i\ge2)\) 已包含在第二个分量中，所以所得分解为 \(V=\bigoplus_i\ker q_i(T)\)。

记 \(V_i=\ker q_i(T)\)，则定义直接给出 \(V_i\subseteq G_{\lambda_i}(T)\)。反过来，取 \(v\in G_{\lambda_i}(T)\)，有某个 \(r\ge1\) 使 \((T-\lambda_iI_V)^rv=0\)。沿上述直和写 \(v=\sum_jv_j\)，其中 \(v_j\in V_j\)。各分量被 \(T\) 保持，故直和的唯一性给出 \((T-\lambda_iI_V)^rv_j=0\)。当 \(j\ne i\) 时，\(q_j\) 与 \((t-\lambda_i)^r\) 互素，Bézout 等式表明一个向量若同时被这两个多项式零化，就只能为零。因此 \(v_j=0\ (j\ne i)\)，于是 \(v\in V_i\)。
\end{proof}

这个分解先将不同特征值的部分分开；在每个 \(G_{\lambda_i}(T)\) 内，\(T-\lambda_iI_V\) 是幂零算子。幂零部分还能怎样分成循环子空间，是 Jordan 块需要进一步描述的问题。

\subsection{线性算子的多项式环作用}
\label{b2:subsec:0103:03}

现在把 \(T\) 连同全部 \(p(T)\) 一起考虑。在 \(V\) 原有的加法上规定
\[
p(t)\cdot v=p(T)v.
\]
算子多项式的加法与乘法公式给出
\[
(p+q)\cdot v=p\cdot v+q\cdot v,\qquad
p\cdot(u+v)=p\cdot u+p\cdot v,\qquad
(pq)\cdot v=p\cdot(q\cdot v),\qquad 1\cdot v=v.
\]
常数多项式的作用与原来的 \(k\) 标量乘法一致。下一章会把这种结构称为 \(k[t]\) 上的模。对于任意维空间，以上定义仍成立；只是算子未必有非零零化多项式。例如在 \(k[t]\) 上令 \(T\) 为乘以 \(t\)，若 \(p(T)=0\)，把它作用在 \(1\) 上即得 \(p(t)=0\)。

\begin{proposition}\label{b2:prop:polynomial-action-subspace}
设 \(V,W\) 是域 \(k\) 上的向量空间，\(U\le V\)，\(T\in\operatorname{End}_k(V)\)、\(S\in\operatorname{End}_k(W)\)，则 \(U\) 被所有 \(p(T)\)（\(p\in k[t]\)）保持，当且仅当它被 \(T\) 保持。对任意 \(k\) 线性映射 \(F:V\to W\)，有
\[
FT=SF
\quad\Longleftrightarrow\quad
F(p(T)v)=p(S)F(v)\quad\text{对所有}\;p\in k[t],v\in V.
\]
\end{proposition}
\begin{proof}
若 \(T(U)\subseteq U\)，反复施加 \(T\) 后仍在 \(U\)，再取有限线性组合就得到 \(p(T)(U)\subseteq U\)；反向取 \(p=t\)。若 \(FT=SF\)，归纳给出 \(FT^j=S^jF\)，包括 \(j=0\)，再由线性性对各项相加即可。反向同样取 \(p=t\)。
\end{proof}

因此，\(T\) 不变子空间恰好是对全部 \(k[t]\) 作用封闭的子空间，下一章称它为子模；关系 \(FT=SF\) 恰好表示 \(F\) 保持 \(k[t]\) 作用，因而是 \(k[t]\) 线性映射。循环子空间 \(k[T]v\) 由一个向量经这种作用生成，称为循环模，而 \(p(T)v=0\) 是这个生成元的关系。这就把不变子空间、算子之间的映射和循环部分的关系放在同一种标量作用下研究。

\ProseParagraphBreak
这种作用也能限制到不变子空间，并下降到商空间。若一个多项式已经零化商算子，它在原空间上的像就落入所取的子空间；再用零化限制算子的多项式，便能零化整个空间。

\begin{proposition}\label{b2:prop:restriction-quotient-polynomials}
设 \(V\) 是域 \(k\) 上的有限维向量空间，\(T\in\operatorname{End}_k(V)\)，\(U\le V\) 满足 \(T(U)\subseteq U\)。则限制 \(T_U=T|_U\) 是 \(U\) 上的算子，且
\[
\bar T:V/U\longrightarrow V/U,\qquad v+U\longmapsto Tv+U
\]
定义了商空间上的算子。商映射 \(q:V\to V/U\) 满足 \(qT=\bar Tq\)，并且
\[
\operatorname{lcm}(\mu_{T_U},\mu_{\bar T})\mid \mu_T
\quad\text{且}\quad \mu_T\mid \mu_{T_U}\mu_{\bar T},
\]
其中最小公倍式取首一者，零空间上的最小多项式约定为 \(1\)。
\end{proposition}
\begin{proof}
若 \(v+U=w+U\)，则 \(v-w\in U\)，不变性给出 \(Tv-Tw\in U\)，所以 \(\bar T\) 与代表元无关。商空间的运算又给出其线性性，且 \(\bar Tq(v)=Tv+U=qT(v)\)。由\cref{b2:prop:polynomial-action-subspace}，对每个 \(p\in k[t]\) 都有
\[
q\,p(T)=p(\bar T)\,q.
\]
若 \(p(T)=0\)，则限制到 \(U\) 后有 \(p(T_U)=0\)，又由上式及 \(q\) 的满射性得到 \(p(\bar T)=0\)。取 \(p=\mu_T\)，\cref{b2:prop:minimal-polynomial-divisibility}给出 \(\mu_{T_U}\mid \mu_T\) 与 \(\mu_{\bar T}\mid \mu_T\)，从而它们的首一最小公倍式也整除 \(\mu_T\)。

令 \(a=\mu_{T_U}\)、\(b=\mu_{\bar T}\)。由于 \(b(\bar T)=0\)，有 \(q\,b(T)=0\)，即 \(b(T)V\subseteq U\)。而 \(a(T)\) 在 \(U\) 上等于 \(a(T_U)=0\)，所以 \(a(T)b(T)=0\)，即 \((ab)(T)=0\)。再用\cref{b2:prop:minimal-polynomial-divisibility}便得 \(\mu_T\mid ab\)。
\end{proof}

乘积一般不能换成最小公倍式。例如在 \(k^2\) 上令 \(Te_1=0\)、\(Te_2=e_1\)，并取 \(U=ke_1\)。限制与商算子都为零，其最小多项式都是 \(t\)；但 \(T\ne0\)、\(T^2=0\)，故 \(\mu_T=t^2\)。子空间和商空间分别只需作用一次就变为零，原空间却可能先从商空间落入子空间，再作用一次才变为零。

% !TeX root = ../../Book2.tex
% 纲要标题：### 1.4 从矩阵到不变量
% 纲要位置：计划纲要.md，第 849 行。

\section{从矩阵到不变量}
\label{b2:sec:0104}

一个算子可以有许多矩阵表示，分类时应先确定哪些变化来自坐标选择。本节写清相似关系，并用小型矩阵说明常见不变量各能辨认什么。最后陈述第四章将要证明的标准形结论，说明为什么需要把算子改写成多项式环的作用。

% 纲要内容（仅作写作计划，尚非教材正文）：
%
% * 换基与相似
% * 有理标准形和 Jordan 标准形的目标
% * 结构定理的证明归于第4章，避免重复证明
%
% ---
%

\subsection{换基与相似}
\label{b2:subsec:0104:01}

设 \(V\) 有限维，有序基为 \(\mathcal B=(e_1,\ldots,e_n)\)。向量 \(v\) 的坐标写为列向量 \([v]_{\mathcal B}\)。若算子 \(T\) 的矩阵为 \(A\)，其第 \(j\) 列就是 \([Te_j]_{\mathcal B}\)，因而
\[
[Tv]_{\mathcal B}=A[v]_{\mathcal B}.
\]
令另一组基 \(\mathcal B'=(e'_1,\ldots,e'_n)\) 满足 \(e'_j=\sum_iP_{ij}e_i\)。由于两组都是基，\(P\) 可逆，且 \([v]_{\mathcal B}=P[v]_{\mathcal B'}\)。于是
\[
[Tv]_{\mathcal B'}=P^{-1}AP[v]_{\mathcal B'}.
\]

\begin{definition}\label{b2:def:matrix-similarity}
设 \(k\) 是域，\(A,B\in M_n(k)\)。若存在 \(P\in\operatorname{GL}_n(k)\) 使 \(B=P^{-1}AP\)，就称 \(A,B\) \term[xiangsi]{相似}[similar]。同一算子在不同有序基下的矩阵恰好构成一个相似类。
\end{definition}

这里左右两侧的变换由同一个换基矩阵联系起来。对于一般线性映射 \(V\to W\)，定义域和陪域可以分别换基。若 \(A,B\in M_{m\times n}(k)\) 满足 \(B=Q^{-1}AP\)，其中 \(P\in\operatorname{GL}_n(k)\)、\(Q\in\operatorname{GL}_m(k)\)，就称 \(A,B\) 矩阵等价。它与算子的相似分类不同：任意初等行、列操作虽然保持矩阵的秩，却未必保持它代表的算子的相似类。

\begin{proposition}\label{b2:prop:similarity-invariants}
设 \(A,B\in M_n(k)\) 是域 \(k\) 上的相似矩阵。它们具有相同的秩、最小多项式、特征多项式及各个 \(p\in k[t]\) 所给的核维数 \(\dim_k\ker p(A)\)。这里把矩阵 \(A\) 的\term[tezheng-duoxiangshi]{特征多项式}[characteristic polynomial]定义为 \(\chi_A(t)=\det(tI_n-A)\)。
\end{proposition}
\begin{proof}
若 \(B=P^{-1}AP\)，归纳得 \(B^j=P^{-1}A^jP\)，故 \(p(B)=P^{-1}p(A)P\)。映射 \(v\mapsto Pv\) 将 \(\ker p(B)\) 同构到 \(\ker p(A)\)，并将 \(\operatorname{im}B\) 同构到 \(\operatorname{im}A\)，所以相应维数不变。又有 \(p(B)=0\iff p(A)=0\)，最小多项式的唯一性给出相等。

在交换环 \(k[t]\) 上使用行列式的乘法公式，得到
\[
\det(tI_n-B)=\det\bigl(P^{-1}(tI_n-A)P\bigr)=\det(tI_n-A).
\]
因此特征多项式也不变。
\end{proof}

这里沿用行列式展开及乘法公式；行列式的外代数解释将在\chapref{b2:ch:09}建立。\(\chi_T\) 因换基不变而也可作为算子本身的记号。但单独知道特征多项式通常不足以确定相似类。例如
\[
A=\begin{pmatrix}0&0\\0&0\end{pmatrix},\qquad
B=\begin{pmatrix}0&1\\0&0\end{pmatrix}
\]
都有特征多项式 \(t^2\)，而 \(m_A=t,m_B=t^2\)，故不相似。秩也相同的例子可以在四维中得到：两个大小为二的零特征值 Jordan 块，与一个大小为三及一个大小为一的块，秩均为二、特征多项式均为 \(t^4\)，最小多项式却分别为 \(t^2,t^3\)。下面会明确这些块的形式。

\subsection{有理标准形与 Jordan 标准形}
\label{b2:subsec:0104:02}

先考察循环向量能提供什么。给定次数 \(d\ge1\) 的首一多项式
\[
p(t)=t^d+a_{d-1}t^{d-1}+\cdots+a_1t+a_0,
\]
其\term[you-juzhen]{友矩阵}[companion matrix]定义为
\[
C(p)=
\begin{pmatrix}
0&0&\cdots&0&-a_0\\
1&0&\cdots&0&-a_1\\
0&1&\cdots&0&-a_2\\
\vdots&\vdots&\ddots&\vdots&\vdots\\
0&0&\cdots&1&-a_{d-1}
\end{pmatrix};
\]
逐列说，\(C(p)e_j=e_{j+1}\ (j<d)\)，而 \(C(p)e_d=-\sum_{i=0}^{d-1}a_ie_{i+1}\)；这也包含 \(d=1\) 时的矩阵 \((-a_0)\)。友矩阵与由代数余子式矩阵转置得到的\term[ban-sui-juzhen]{伴随矩阵}[adjugate matrix]是不同对象；本书分别使用这两个名称。

\ProseParagraphBreak
若 \(v\) 为循环向量，且 \(m_T=p\)，则\cref{b2:prop:cyclic-basis}给出基 \(v,Tv,\ldots,T^{d-1}v\)。前 \(d-1\) 个基向量各被送到后一个，而
\[
T^dv=-a_0v-a_1Tv-\cdots-a_{d-1}T^{d-1}v,
\]
所以算子在这组基下的矩阵就是 \(C(p)\)。这已完整解释单个循环子空间的矩阵形式。一般算子怎样分成循环部分，以及不同分法为什么给出相同的分类数据，仍需要证明。

\begin{claim}[有理标准形，证明见第四章]\label{b2:claim:rational-canonical-goal}
设 \(V\ne0\) 是域 \(k\) 上的有限维向量空间，\(T\in\operatorname{End}_k(V)\)。在适当的有序基下，\(T\) 的矩阵为
\[
\operatorname{diag}\bigl(C(d_1),\ldots,C(d_r)\bigr),
\qquad d_1\mid d_2\mid\cdots\mid d_r,
\]
其中 \(d_i\in k[t]\) 首一且次数为正，\(\sum_i\deg d_i=\dim_k V\)。这些多项式唯一，且 \(m_T=d_r\)、\(\chi_T=\prod_i d_i\)。
\end{claim}

这一形式称为\term[youli-biaozhunxing]{有理标准形}[rational canonical form]。“有理”不要求 \(k=\mathbb Q\)，而是说所有构造都在原基域上进行，不要求先把多项式分解成一次因子。存在性与唯一性将在\cref{b2:thm:rational-canonical-form}证明，最小多项式与特征多项式公式见\cref{b2:prop:operator-characteristic-minimal}；本章的证明与习题均不依赖这一前瞻陈述。

\ProseParagraphBreak
若多项式在 \(k\) 上分裂，还可使用更细的块。定义
\[
J_d(\lambda)=
\begin{pmatrix}
\lambda&1&&0\\
&\lambda&\ddots&\\
&&\ddots&1\\
0&&&\lambda
\end{pmatrix},\qquad J_1(\lambda)=(\lambda).
\]
它称为\term[Jordan-kuai]{Jordan 块}[Jordan block]。减去 \(\lambda I\) 后，每个 \(e_j\ (j>1)\) 被送到 \(e_{j-1}\)，\(e_1\) 被送到零；因而第 \(r\) 次幂的核维数为 \(\min(r,d)\)。

\begin{claim}[Jordan 标准形，证明见第四章]\label{b2:claim:jordan-canonical-goal}
设 \(V\) 是域 \(k\) 上的有限维向量空间，\(T\in\operatorname{End}_k(V)\)。若 \(T\) 的特征多项式在 \(k\) 上分裂为一次因子的乘积，则在适当的有序基下，\(T\) 的矩阵是若干 Jordan 块 \(J_d(\lambda)\)（\(d\ge1,\lambda\in k\)）的直和。这些块除排列次序外唯一。
\end{claim}

这称为\term[Jordan-biaozhunxing]{Jordan 标准形}[Jordan canonical form]。完整证明归于\cref{b2:thm:jordan-canonical-form}。核维数序列说明了块大小为什么可望成为可恢复的不变量：对固定特征值，一个大小为 \(d\) 的块在第 \(r\) 次幂的核中贡献 \(\min(r,d)\) 维；相邻差便计算大小至少为 \(r\) 的块数。这是分析已给定块矩阵的直接计算，还不是每个算子都有这种分解的证明。

\subsection{标准形的模论解释}
\label{b2:subsec:0104:03}

对于单个循环子空间，多项式与向量之间有一个明确的对应。

\begin{proposition}\label{b2:prop:cyclic-polynomial-quotient}
设 \(V\ne0\) 是域 \(k\) 上的有限维向量空间，\(T\in\operatorname{End}_k(V)\)，\(0\ne v\in V\)，且 \(p=m_{T,v}\) 是零化 \(v\) 的首一最小多项式。映射
\[
k[t]/(p)\longrightarrow k[T]v,\qquad f+(p)\longmapsto f(T)v
\]
是 \(k\) 线性同构，并把左侧乘以 \(t\) 的运算变为右侧施加 \(T\) 的运算。
\end{proposition}
\begin{proof}
映射 \(k[t]\to k[T]v\)、\(f\mapsto f(T)v\) 线性且满，其核恰好是 \(p\) 的全部倍数，即理想 \((p)\)。\cref{b2:prop:linear-first-isomorphism}给出商空间到像的线性同构。对于任何 \(f\)，\((tf)(T)v=T(f(T)v)\)，证明最后的相容性。
\end{proof}

商环 \(k[t]/(p)\) 在这里同时有向量空间结构和多项式作用。带余除法说明 \(1,t,\ldots,t^{d-1}\) 的陪集是基；用这组基表示乘 \(t\)，又得到 \(C(p)\)。后续的模结构定理要做的，就是把一般有限维算子分解为这样的循环商，并识别每一部分的关系多项式。本节只证明了单个循环部分的对应，不以标准形的存在性作为它的依据。

\ProseParagraphBreak
基域对分类形式的影响可以从一个二维例子看出。在 \(\mathbb R^2\) 上取
\[
A=\begin{pmatrix}0&-1\\1&0\end{pmatrix}=C(t^2+1).
\]
因为 \(A^2=-I\)，且 \(e_1,Ae_1=e_2\) 是基，故最小多项式为 \(t^2+1\)。这个多项式在 \(\mathbb R\) 上没有根，算子没有实特征向量。把同一矩阵放在 \(\mathbb C^2\) 上，向量 \((1,-i)\)、\((1,i)\) 分别对应特征值 \(i,-i\)，而且二者独立，所以矩阵在复数域上相似于 \(\operatorname{diag}(i,-i)\)。矩阵条目没有变化，可用的标量与基向量却变了。讨论标准形时，必须同时写明所工作的域。


\begin{chaptersummary}
基提供坐标，商空间和映射的核、像则可以直接由对象定义。商空间及核、像的维数加法公式在任意维数下都成立；有限维时才能用整数减法恢复商或像的维数。无限基数不能这样相减或消去，无限维空间的单射自映射与满射自映射也不再等同。

\ProseParagraphBreak
对偶把向量的研究转成线性泛函的研究。对偶映射反转箭头，求值映射则不需选基；有限维空间与其双对偶自然同构，而任意选基给出的空间与对偶之间的同构带有额外选择。无限直和的对偶是直积，无限维求值映射也可能不满射。

\ProseParagraphBreak
线性算子使多项式环作用在空间上。循环子空间由一个向量及其关系多项式描述，友矩阵记录这组关系；相似变换保持这种作用的结构。当最小多项式在基域中分裂时，互素核分解把空间分成各广义特征子空间，在对应于 \(\lambda\) 的部分上，\(T-\lambda I\) 是幂零算子。将这个幂零作用分成循环链，就得到 Jordan 块；每条链的长度由零化它的 \((t-\lambda)\) 幂记录。一般循环分解的存在性和唯一分类数据，以及这些链能否组成全空间的基，将在第四章用模的结构定理解决。下一章将向量空间的标量作用推广到一般环上的模；非零标量不再必然可逆，基、补空间及相关分裂结论都须重新检查。
\end{chaptersummary}

\BookExercises

\begin{exercise}\label{b2:ex:01-evaluation-quotient}
在 \(V=k[t]_{\le3}\) 中，令 \(U=\{p\in V\mid p(0)=p(1)=0\}\)。求 \(U\) 的一组基及一个补空间，构造 \(V/U\cong k^2\) 的显式同构。这里 \(k\) 可以有任意特征。
\end{exercise}
\begin{solution}
\(t\) 与 \(t-1\) 互素，所以两个取值为零当且仅当 \(t(t-1)\mid p\)。故 \(U\) 的基为 \(t(t-1),t^2(t-1)\)。令 \(W=\operatorname{span}_k(1,t)\)。对任意 \(p\in V\)，多项式
\[
r(t)=(1-t)p(0)+t\cdot p(1)
\]
在 \(0,1\) 处与 \(p\) 同值，因而 \(p-r\in U\)；\(W\) 中同时在这两点为零的多项式只能为零。因此 \(V=U\oplus W\)。映射 \(p\mapsto(p(0),p(1))\) 的核为 \(U\)，并由上述插值式知它满射。由\cref{b2:prop:linear-first-isomorphism}，所求同构为 \(p+U\mapsto(p(0),p(1))\)，其逆把 \((a,b)\) 送到 \((1-t)a+tb+U\)。域中 \(0\ne1\)，所以任意特征均适用。
\end{solution}

\begin{exercise}\label{b2:ex:01-quotient-map-test}
设 \(V=k^2\)、\(U=k(e_1+e_2)\)，线性映射 \(f,g:V\to k\) 分别由
\[
f(e_1)=1,\quad f(e_2)=-1,\qquad
g(e_1)=1,\quad g(e_2)=0
\]
确定。对每个映射，判断它能否写成 \(V\xrightarrow{q}V/U\xrightarrow{h}k\)，其中 \(q\) 是商映射，\(h\) 为线性映射；若能，写出 \(h\)，若不能，说明原因。
\end{exercise}
\begin{solution}
因为 \(f(e_1+e_2)=0\)，有 \(U\subseteq\ker f\)。按\cref{b2:prop:quotient-linear-universal}，存在唯一的线性映射 \(\bar f:V/U\to k\) 使 \(f=\bar f q\)，其公式为 \(\bar f((x,y)+U)=x-y\)。若换用同一陪集的另一个代表元 \((x+c,y+c)\)，所得值仍是 \(x-y\)，所以也可直接核对良定义。另一方面，\(g(e_1+e_2)=1\)，而 \(q(e_1+e_2)=0\)；若 \(g=\bar gq\)，便会有 \(1=\bar g(0)=0\)，矛盾。
\end{solution}

\begin{exercise}\label{b2:ex:01-complement-graphs}
设 \(V=U\oplus W\)。证明 \(U\) 的全部补空间恰为
\[
W_h=\{h(w)+w\mid w\in W\},\qquad h\in\operatorname{Hom}_k(W,U),
\]
并且不同 \(h\) 给出不同补空间。
\end{exercise}
\begin{solution}
\(W_h\) 对线性组合封闭。若 \(h(w)+w\in U\)，沿已给直和取 \(W\) 分量即得 \(w=0\)，故交为零。又有 \(u+w=(u-h(w))+(h(w)+w)\)，所以 \(V=U\oplus W_h\)。

反过来，设 \(W'\) 为补空间。原分解所给的投影 \(\pi_W:V\to W\) 在 \(W'\) 上单射，因为其核为 \(W'\cap U=0\)；它也满射，因为每个 \(w\in W\) 可以沿 \(U\oplus W'\) 分解。其逆线性，将 \(w\) 写成唯一的 \(h(w)+w\)，其中 \(h(w)\in U\)。取 \(U\) 分量可知 \(h\) 线性，且由 \(W'\) 唯一确定。
\end{solution}

\begin{exercise}\label{b2:ex:01-shift-dual}
令 \(V=k^{(\mathbb N)}\)，基为 \(e_1,e_2,\ldots\)。定义 \(S(e_i)=e_{i+1}\)，\(L(e_1)=0\)、\(L(e_{i+1})=e_i\)。计算 \(LS,SL\)，并分别判断 \(S,L\) 的单射性和满射性；再在 \(V^*\cong k^{\mathbb N}\) 上写出 \(S^*,L^*\) 的公式及其核与像。
\end{exercise}
\begin{solution}
在基向量上计算得 \(LS=I_V\)，而 \(SL\) 将 \(e_1\) 送到零并固定其余基向量。\(S\) 单射，像为第一坐标为零的有限支集族；\(L\) 满射，核为 \(ke_1\)。按\cref{b2:prop:dual-direct-sum}，泛函由任意族 \(a_i=\varphi(e_i)\) 表示。于是
\[
S^*(a_1,a_2,\ldots)=(a_2,a_3,\ldots),\qquad
L^*(a_1,a_2,\ldots)=(0,a_1,a_2,\ldots).
\]
\(S^*\) 满射，核为 \(k(1,0,\ldots)\)；\(L^*\) 单射，像是第一坐标为零的全部序列。对偶映射反转复合次序：\(S^*L^*=I_{V^*}\)，与 \((LS)^*=I_{V^*}\) 一致。
\end{solution}

\begin{exercise}\label{b2:ex:01-double-dual-concrete}
令 \(V=k^2\)、\(W=k^3\)，并定义 \(f:V\to W\) 为 \(f(x,y)=(x+y,y,x)\)。分别写出 \(f^*:W^*\to V^*\) 与 \(f^{**}:V^{**}\to W^{**}\) 的作用，并对任意 \((x,y)\in V\) 和 \(\lambda\in W^*\) 直接验证 \(j_W(f(x,y))(\lambda)=f^{**}(j_V(x,y))(\lambda)\)。
\end{exercise}
\begin{solution}
用标准基的对偶基分别表示 \(V^*\) 与 \(W^*\)，把 \(\lambda\in W^*\) 记作 \((a,b,c)\)，即 \(\lambda(u,v,w)=au+bv+cw\)。则
\[
f^*(a,b,c)=(a+c,a+b).
\]
对 \(F\in V^{**}\)，双对偶映射按定义满足 \(f^{**}(F)(\lambda)=F(f^*(\lambda))\)。在 \(V^*\) 的上述对偶基所对应的双对偶基下，\(F=(u,v)\) 表示 \(F(s,t)=us+vt\)。于是 \(f^{**}(F)(a,b,c)=u(a+c)+v(a+b)=(u+v)a+vb+uc\)，按 \(W^*\) 的对偶基取值读出 \(f^{**}(u,v)=(u+v,v,u)\)。特别地，\(j_V(x,y)\) 在坐标为 \((s,t)\) 的泛函上取值 \(sx+ty\)，所以
\[
f^{**}(j_V(x,y))(\lambda)=(a+c)x+(a+b)y
=a(x+y)+by+cx=j_W(f(x,y))(\lambda).
\]
该计算不要求 \(k\) 的特征为零。
\end{solution}

\begin{exercise}\label{b2:ex:01-trace-pairing}
设 \(n\ge1\)，在 \(M_n(k)\) 上定义双线性配对 \(\beta(A,B)=\operatorname{tr}(AB)\)。证明它在任意特征下非退化。令所有上三角矩阵组成的子空间为
\[
U=\{A\in M_n(k):A_{ij}=0\;\text{当}\;i>j\},
\]
并定义 \(U^{\perp_\beta}=\{B\in M_n(k):\beta(A,B)=0\;\text{对所有}\;A\in U\}\)。求 \(U^{\perp_\beta}\)，并解释为什么 \(\operatorname{char}k\mid n\) 时 \(\operatorname{tr}(I_n)=0\) 不与非退化性矛盾。
\end{exercise}
\begin{solution}
设 \(E_{ij}\) 为矩阵单位。展开迹得 \(\operatorname{tr}(AB)=\sum_{i,j}A_{ij}B_{ji}\)，所以 \(\beta(A,E_{ji})=A_{ij}\)。若 \(A\ne0\)，选一个非零条目即可找到使配对非零的 \(B\)；交换左右同理。上三角空间由 \(E_{ij}\ (i\le j)\) 生成，条件 \(\beta(E_{ij},B)=B_{ji}=0\) 排除 \(B\) 的对角线及以下条目，所以 \(U^{\perp_\beta}\) 是严格上三角空间。

\(\operatorname{tr}(I_n)=\beta(I_n,I_n)=0\) 只说明 \(I_n\) 与自身的配对为零。非退化要求一个非零矩阵不与所有矩阵都配成零；例如 \(\beta(I_n,E_{11})=1\)。
\end{solution}

\begin{exercise}\label{b2:ex:01-basis-pairing}
在 \(V=\mathbb Q^2\) 的标准基 \(\mathcal B\) 下，用 \(\beta(v,w)=[v]_{\mathcal B}^{\mathsf T}[w]_{\mathcal B}\) 定义配对。若 \(\mathcal B'\) 的换基矩阵为 \(P\)，求该配对在新坐标下的矩阵。取 \(P=\begin{pmatrix}1&1\\0&1\end{pmatrix}\)，说明“把一组基送到其对偶基”给出的 \(V\to V^*\) 同构为何依赖所选基。
\end{exercise}
\begin{solution}
由 \([v]_{\mathcal B}=P[v]_{\mathcal B'}\)，有
\[
\beta(v,w)=[v]_{\mathcal B'}^{\mathsf T}P^{\mathsf T}P[w]_{\mathcal B'}.
\]
所给 \(P\) 满足 \(P^{\mathsf T}P=\begin{pmatrix}1&1\\1&2\end{pmatrix}\ne I\)。定义域取 \(\mathcal B'\)、陪域取其对偶基时，原同构 \(v\mapsto\beta(v,-)\) 的矩阵为 \(P^{\mathsf T}P\)；把 \(\mathcal B'\) 的每个基向量直接送到相应的对偶基，所得同构的矩阵则是 \(I\)。二者不同。双对偶中的求值映射 \(j_V(v)(\varphi)=\varphi(v)\) 则根本没有选基步骤，不能把这两种同构混为一谈。
\end{solution}

\begin{exercise}\label{b2:ex:01-infinite-annihilator}
仍取 \(V=k^{(\mathbb N)}\)，记坐标泛函为 \(\varepsilon_i\)，令 \(E\subseteq V^*\) 由 \(\varepsilon_i-\varepsilon_{i+1}\ (i\ge1)\) 生成。在 \(V^*\cong k^{\mathbb N}\) 的识别下，刻画 \(E\) 所对应的数列；证明 \(E\ne V^*\)，而被 \(E\) 的全部泛函同时零化的向量只有零。
\end{exercise}
\begin{solution}
\(E\) 恰由有限支集且坐标和为零的泛函组成。每个生成元都如此；反过来，若 \(\sum_{i=1}^ma_i=0\)，取 \(b_j=a_1+\cdots+a_j\)，则
\[
\sum_{i=1}^ma_i\varepsilon_i
=\sum_{j=1}^{m-1}b_j(\varepsilon_j-\varepsilon_{j+1}).
\]
因此 \(\varepsilon_1\notin E\)，所以 \(E\) 为真子空间。若 \(v=(v_i)\) 被所有生成元零化，就有 \(v_i=v_{i+1}\)，故坐标全相等；有限支集又迫使它们全为零。这说明在无限维中，由 \(V\) 中向量形成的求值泛函，未必足以把 \(V^*\) 的每个真子空间与一个外部元素分开。
\end{solution}

\begin{exercise}\label{b2:ex:01-derivative-dual}
令 \(V=k[t]_{\le3}\)，\(D\) 为形式求导算子。分别在 \(\operatorname{char}k=0\) 和 \(\operatorname{char}k=2\) 时求 \(\ker D\)、\(\operatorname{im}D\)、\(\mu_D\)，并判断是否有循环向量。若 \(\varepsilon_j\) 取多项式的 \(t^j\) 系数，写出 \(D^*\varepsilon_j\)。
\end{exercise}
\begin{solution}
由 \(D(t^j)=jt^{j-1}\)，特征零时核为 \(k\)，像为 \(k[t]_{\le2}\)。\(D^4=0\)，而 \(D^3(t^3)=6\ne0\)，故 \(\mu_D=t^4\)；\(t^3,3t^2,6t,6\) 是基，所以 \(t^3\) 为循环向量。

特征二时 \(D(1)=D(t^2)=0\)、\(D(t)=1\)、\(D(t^3)=t^2\)，故核与像都为 \(\operatorname{span}_k(1,t^2)\)。此时 \(D^2=0\) 且 \(D\ne0\)，所以 \(\mu_D=t^2\)，任一循环子空间维数至多二，不能等于四维的 \(V\)。两种特征下统一有
\[
D^*\varepsilon_j=(j+1)\varepsilon_{j+1}\quad(0\le j<3),\qquad
D^*\varepsilon_3=0,
\]
其中整数系数取其在 \(k\) 中的像。
\end{solution}

\begin{exercise}\label{b2:ex:01-cyclic-characteristic}
在 \(k^3\) 上定义 \(Te_1=e_2,Te_2=e_3,Te_3=2e_1\)。分别取 \(k=\mathbb Q\) 与 \(k=\mathbb F_3\)，求 \(\mu_T\)，并判断 \(v=e_1+e_2\) 是否为循环向量。
\end{exercise}
\begin{solution}
在两个域上，\(e_1,Te_1,T^2e_1\) 都是标准基，而 \(T^3=2I\)，所以\cref{b2:prop:cyclic-basis}给出 \(\mu_T=t^3-2\)。以 \(v,Tv,T^2v\) 为列所得矩阵是
\[
P=\begin{pmatrix}1&0&2\\1&1&0\\0&1&1\end{pmatrix},\qquad\det P=3.
\]
故在 \(\mathbb Q\) 上 \(v\) 循环。在 \(\mathbb F_3\) 上，\(\mu_T=(t+1)^3\)，且 \(v=(T+I)e_1\)，所以 \((T+I)^2v=0\)。又有 \((T+I)v=e_1+2e_2+e_3\ne0\)，故 \(\mu_{T,v}=(t+1)^2\)。其循环子空间只有二维，\(v\) 不循环。
\end{solution}

\begin{exercise}\label{b2:ex:01-invariant-subspaces}
设 \(T\) 在 \(k^3\) 上的矩阵为 \(\operatorname{diag}(\lambda,\lambda,\mu)\)，其中 \(\lambda\ne\mu\)。证明每个 \(T\) 不变子空间都形如 \(U_\lambda\oplus U_\mu\)，其中 \(U_\lambda\le\operatorname{span}_k(e_1,e_2)\)，而 \(U_\mu\) 为 \(0\) 或 \(ke_3\)。并求最小多项式，判断是否存在循环向量。
\end{exercise}
\begin{solution}
算子 \(P_\lambda=(T-\mu I)/(\lambda-\mu)\) 和 \(P_\mu=(T-\lambda I)/(\mu-\lambda)\) 分别是到两个特征子空间的投影，且和为 \(I\)。若 \(U\) 被 \(T\) 保持，由\cref{b2:prop:polynomial-action-subspace}也被这两个多项式保持。因此每个 \(u\in U\) 的两个特征分量仍在 \(U\)，给出所述分解；反向封闭性直接成立。\(U_\mu\) 所在空间一维，所以只有两种选择。

多项式 \(p\) 零化 \(T\) 当且仅当 \(p(\lambda)=p(\mu)=0\)，故 \(\mu_T=(t-\lambda)(t-\mu)\)。若存在循环向量，由\cref{b2:prop:cyclic-basis}其次数应为三，与实际次数二矛盾。
\end{solution}

\begin{exercise}\label{b2:ex:01-matrix-equivalence-rank}
证明同尺寸的两个域上矩阵等价，当且仅当它们的秩相同。再取
\[
A=\begin{pmatrix}1&0\\0&0\end{pmatrix},\qquad
B=\begin{pmatrix}0&1\\0&0\end{pmatrix}.
\]
给出从 \(A\) 到 \(B\) 的初等行或列操作，并说明两者为什么不相似。
\end{exercise}
\begin{solution}
若 \(B=Q^{-1}AP\)，其中 \(P,Q\) 可逆，则左右乘可逆矩阵不改变线性映射的秩。反过来，设 \(A\in M_{m\times n}(k)\) 的秩为 \(r\)。先取 \(\ker A\) 的一组基，再添入 \(v_1,\ldots,v_r\) 扩充为定义域的基；\(Av_1,\ldots,Av_r\) 构成 \(\operatorname{im}A\) 的基，将它们扩充为陪域的基。定义域以 \(v_1,\ldots,v_r\) 排在核的基向量之前，陪域以 \(Av_1,\ldots,Av_r\) 排在其余基向量之前，则 \(A\) 的矩阵成为左上角是 \(I_r\)、其余位置是零的矩形矩阵 \(D_r\)。对 \(B\) 作同样处理；若秩也为 \(r\)，两者都等价于 \(D_r\)，故互相等价。

将 \(A\) 的两列交换便得到 \(B\)。然而 \(A^2=A\)，而 \(B^2=0\ne B\)。相似矩阵满足相同的多项式等式，因此 \(A,B\) 不相似。这一反例在任意域上成立。
\end{solution}

\begin{exercise}\label{b2:ex:01-insufficient-invariants}
在已经写成若干 \(J_d(0)\) 直和的矩阵中，求这样的最小阶数 \(n\)：存在两个 \(n\) 阶矩阵，秩、特征多项式与最小多项式都相同，却不相似。给出达到该阶数的例子，并证明更小的阶数不可能。这里不要求先证明一般幂零算子具有 Jordan 分解。
\end{exercise}
\begin{solution}
若块大小为 \(d_1,\ldots,d_r\)，则矩阵阶数为 \(\sum_i d_i=n\)，秩为 \(n-r\)，特征多项式为 \(t^n\)，最小多项式为 \(t^{\max_i d_i}\)。最后一个公式可从循环块的关系直接看出：零化 \(J_d(0)\) 的多项式恰为 \(t^d\) 的倍数。因此题目要求块大小的两个不同分组具有相同的总和、项数和最大项。

阶数七时，可以取
\[
A=J_3(0)\oplus J_3(0)\oplus J_1(0),\qquad
B=J_3(0)\oplus J_2(0)\oplus J_2(0).
\]
它们秩均为四，特征多项式均为 \(t^7\)，最小多项式均为 \(t^3\)，但平方的核维数分别为五与六，所以不相似。

以下设 \(n\le6\)，固定共同的块数 \(r\) 和最大块大小 \(D\)。若 \(D\le2\)，设大小二的块有 \(s\) 个，则 \(n=r+s\)，已经唯一确定这组块。若 \(D\ge4\)，取走一个大小 \(D\) 的块后，剩余大小之和至多为二；固定剩余块数 \(r-1\) 后，剩余分组也唯一。最后设 \(D=3\)：取走一个三阶块后，剩余大小之和至多为三。对和为零、一、二、三的分组，在项数固定时都唯一，分别只需考虑空分组、\(1\)、\(2\) 或 \(1+1\)、\(3\) 或 \(2+1\) 或 \(1+1+1\)。所以给定 \(n,r,D\) 后，块大小已唯一；两个矩阵至多只差块的排列，换基即可相互得到。故最小阶数为七。
\end{solution}

\begin{exercise}\label{b2:ex:01-cyclic-centralizer}
若非零有限维空间 \(V\) 上的算子 \(T\) 有循环向量，证明与 \(T\) 交换的全部线性算子恰为 \(k[T]=\{p(T)\mid p\in k[t]\}\)，并求这个空间的维数。
\end{exercise}
\begin{solution}
取循环向量 \(v\)。若 \(FT=TF\)，可写 \(F(v)=p(T)v\)。对任意 \(q(T)v\)，有
\[
F(q(T)v)=q(T)F(v)=q(T)p(T)v=p(T)q(T)v.
\]
这些向量遍及 \(V\)，故 \(F=p(T)\)。反之，所有算子多项式都与 \(T\) 交换。令 \(n=\dim V=\deg \mu_T\)。带余除法说明 \(I,T,\ldots,T^{n-1}\) 生成 \(k[T]\)，最小多项式的定义说明它们线性无关，所以维数为 \(n\)。
\end{solution}

\begin{exercise}\label{b2:ex:01-restriction-quotient-polynomials}
设 \(V\) 为域 \(k\) 上的有限维向量空间，\(0\ne U\subsetneq V\) 是算子 \(T\) 的不变子空间。限制算子 \(T_U\) 与商算子 \(\bar T\) 都幂零，其幂零指数分别为 \(a,b\ge1\)，即零化它们的最小正幂分别为 \(a,b\)。证明 \(T\) 幂零且其指数 \(c\) 满足
\[
\max(a,b)\le c\le a+b.
\]
对任意 \(a,b\ge1\)，构造达到上界的例子。若 \(U\) 有 \(T\) 不变的补空间，再证明 \(c=\max(a,b)\)，并给出达到此值的例子。
\end{exercise}
\begin{solution}
此时 \(\mu_{T_U}=t^a\)、\(\mu_{\bar T}=t^b\)。由\cref{b2:prop:restriction-quotient-polynomials}，
\[
t^{\max(a,b)}\mid \mu_T\mid t^{a+b}.
\]
因此 \(\mu_T=t^c\) 且 \(\max(a,b)\le c\le a+b\)，这也说明 \(T\) 确实幂零。

取 \(T=J_{a+b}(0)\) 作用在 \(k^{a+b}\) 上，令 \(U=\operatorname{span}_k(e_1,\ldots,e_a)\)。因 \(Te_j=e_{j-1}\)（\(j>1\)）、\(Te_1=0\)，\(U\) 不变且限制为 \(J_a(0)\)；商空间以 \(e_{a+1}+U,\ldots,e_{a+b}+U\) 为基，商算子为 \(J_b(0)\)。原算子的幂零指数为 \(a+b\)，达到上界。

若 \(V=U\oplus W\) 且 \(W\) 也不变，商映射在 \(W\) 上给出 \(W\cong V/U\)，并将 \(T|_W\) 对应到 \(\bar T\)。故 \(T^r=0\) 当且仅当 \(T_U^r=0\) 和 \(\bar T^r=0\) 同时成立，最小的这样的 \(r\) 是 \(\max(a,b)\)。具体取 \(T=J_a(0)\oplus J_b(0)\)，并令 \(U\) 为第一个块所在的空间，即达到这个值。上界例子把两层接成一条长链，直和例子则把两层分开。
\end{solution}
